\documentclass[10pt,a4paper]{article}

% ----------------- Paquetes -------------------%
\usepackage[T1]{fontenc}
\usepackage[utf8]{inputenc}
\usepackage[a4paper,margin=2.5cm]{geometry}
\usepackage{mathptmx}             
\usepackage{changepage} 
\usepackage{setspace}
\usepackage[spanish]{babel}
\usepackage[labelfont=bf,labelsep=space]{caption}
\usepackage{amssymb}
\usepackage{amsmath}
\usepackage{ebgaramond}
\usepackage{placeins}
\usepackage{etoolbox}
\usepackage{setspace}
\usepackage{parskip} 
\usepackage{tcolorbox}
\usepackage{needspace}
\usepackage{xparse}
\usepackage{ocgx2}
\usepackage{hyperref}
\usepackage{hypcap}
\usepackage{multirow} 



%%%%%%%%%%%%%%%%%%%%%%%%%%%%%%%%%%%%%%%%%%%%%%%%%%%%%%%%%%%%%%%%
% --- Fija primero tus valores globales "buenos" ---
\setstretch{1.2}                 % Interlineado global
\setlength{\parskip}{0.7\baselineskip}
\setlength{\parindent}{0pt}
\newlength{\GlobalParskip}
\setlength{\GlobalParskip}{\parskip}

\newcounter{eqnum}[subsection]
\renewcommand{\theeqnum}{\arabic{eqnum}}
\newcounter{alphaitem}
\newcounter{numitem}

%% ------------------- Recuadros----------------------%%
\newcommand{\puntosclave}[1]{%
  \begin{tcolorbox}[
    colframe=black,
    title={Puntos clave},
    before upper={\setlength{\parindent}{0.5cm}\setlength{\parskip}{0.5\baselineskip}}
  ]
  #1
  \end{tcolorbox}
}

\newcommand{\modificaciones}[1]{
  \begin{tcolorbox}[
    colback=white,
    colframe=red,
    title={Modificaciones a realizar},
    before upper={\setlength{\parindent}{0.5cm}\setlength{\parskip}{0.5\baselineskip}}
  ]
  #1
  \end{tcolorbox}
}

%% ------------------- Preguntas TEST----------------------%%
% Opciones: radio (single-choice)
\NewDocumentCommand{\OptRadio}{m m m}{%
  \noindent\CheckBox[radio,name=#1,value=#2,width=1.2ex,height=1.2ex]{}%
  \;\textbf{#2.}~#3\par
}

% Opciones: checkbox (multi-choice)
\NewDocumentCommand{\OptCheck}{m m}{%
  \noindent\CheckBox[name=#1,width=1.2ex,height=1.2ex,bordercolor={0 0 0}]{}%
  \;#2\par
}

% Pregunta single-choice (radio)
% Args: 1=id, 2=enunciado, 3=opciones, 4=solución+justificación, 5=imagen (o {}), 6=origen
\NewDocumentCommand{\PreguntaTestSingle}{m m m m m m}{%
  \par\medskip
  \noindent\textbf{#1.}~#2\par\smallskip

    % Imagen (si no es {})
    \IfBlankTF{#5}{}{%
      \begin{center}
        \includegraphics[width=0.75\linewidth]{#5}
      \end{center}
    }

    \begin{Form}
      #3
    \end{Form}

    \par\smallskip
    \switchocg{sol-#1}{\fbox{Mostrar / ocultar solución}}%
    \hfill{\footnotesize #6}\par\smallskip

    \begin{ocg}{Solución}{sol-#1}{0}
      \noindent #4
    \end{ocg}
  \par\medskip
}

% Pregunta multi-choice (checkboxes)
\NewDocumentCommand{\PreguntaTestMulti}{m m m m m m}{%
  \par\medskip
  \noindent\textbf{#1.}~#2\par\smallskip

    % Imagen (si no es {})
    \IfBlankTF{#5}{}{%
      \begin{center}
        \includegraphics[width=0.75\linewidth]{#5}
      \end{center}
    }
    \begin{Form}
      #3
    \end{Form}

    \par\smallskip
    \switchocg{sol-#1}{\fbox{Mostrar / ocultar solución}}%
    \hfill{\footnotesize #6}\par\smallskip

    \begin{ocg}{Solución}{sol-#1}{0}
      \noindent #4
    \end{ocg}
  \par\medskip
}

%% ------------------- Listas----------------------%%
    % --- Lista alfabética ---
    \newenvironment{la}
      {\begin{list}{\alph{alphaitem})}{%
          \usecounter{alphaitem}%
          \setlength{\leftmargin}{1cm}%
          \setlength{\parskip}{3pt}% 
          \setlength{\itemsep}{0pt}% ← espacio entre ítems
          \setlength{\parsep}{0pt}%  ← párrafos dentro de un ítem
      }}
      {\end{list}}
      
    % --- Lista numérica ---
    \newenvironment{lnum}
      {\begin{list}{\arabic{numitem}.}{%
          \usecounter{numitem}%
          \setlength{\leftmargin}{1cm}%
          \setlength{\parskip}{3pt}% 
          \setlength{\itemsep}{0pt}% ← espacio entre ítems
          \setlength{\parsep}{0pt}%  ← párrafos dentro de un ítem
      }}
      {\end{list}}

%% ------------------- Preguntas Test ----------------------%%
      \usepackage{xparse} % si ya lo tienes, no lo repitas

    \NewDocumentEnvironment{test}{m m o}{%
      \begin{tcolorbox}[
        colback=white,
        colframe=black!15,
        boxrule=0.6pt,
        arc=2mm,
        left=2mm,right=2mm,top=2mm,bottom=2mm
      ]
      \centering
      \includegraphics[width=\linewidth]{#1}
      \end{tcolorbox}
    
      \vspace{2mm}
    
      % Caja SOLO para texto (SÍ partible y mantiene márgenes al partir)
      \begin{tcolorbox}[
        enhanced,
        breakable,
        colback=black!3,
        colframe=black!25,
        boxrule=0.6pt,
        arc=2mm,
        left=2mm,right=2mm,top=1.5mm,bottom=1.5mm
      ]
      \textbf{Respuesta:} \textbf{#2}%
      \IfValueT{#3}{\par\smallskip \textbf{Justificación:} #3}
      \end{tcolorbox}
    }{}
      
% ------------------- Imágenes----------------------%
    \graphicspath{{Img/}}% Rutas por defecto 
    % --- Imagen adaptativa ---
    % Registros de dimensión definidos una sola vez
    \newdimen\imgcap
    \newdimen\imgmin
    \newdimen\imgmax
    \newdimen\imgremain
    \newdimen\imgh
    \newdimen\imgneed
    
    \newcommand{\imagenfit}[3]{%
      \addvspace{10pt}% separación antes
      \par\begingroup
        % Parámetros de composición (asignaciones, no \newdimen)
        \imgcap=1\baselineskip            % alto estimado de título+fuente
        \imgmin=0.15\textheight
        \imgmax=0.15\textheight
        \imgremain=\pagegoal
        \ifdim\pagegoal=\maxdimen \imgremain=0pt\fi
        \advance\imgremain by -\pagetotal
        \advance\imgremain by -\imgcap
        % Decisión de altura destino
        \ifdim\imgremain<\imgmin
          \newpage
          \imgh=0.20\textheight
        \else
          \imgh=\imgremain
          \ifdim\imgh>\imgmax \imgh=\imgmax \fi
        \fi
        % Reservar espacio para evitar cortes del bloque completo
        \imgneed=\imgh \advance\imgneed by \imgcap
        \Needspace{\imgneed}
        % Cuerpo
        \noindent\begin{minipage}{\textwidth}
          \centering
          \captionof{figure}{\textemdash\ #2}\par\vspace{0.3em}% título arriba
          \includegraphics[width=\linewidth,height=\imgh,keepaspectratio]{\detokenize{#1}}\par
          \vspace{0.3em}\small\textit{Fuente: }#3% fuente abajo
        \end{minipage}%
      \endgroup
      \par\addvspace{10pt}%
    }

%------------ Bloque de RECUADRO ESPECIAL -------------%
    \newenvironment{specialblock}[1]{%
      \begin{quote} % crea sangría a izquierda y derecha
      \small % texto más pequeño
      \noindent\textbf{#1}\\[-0.3em] % título en negrita
      \rule{\linewidth}{0.4pt}\\[0.6em]}{
      \end{quote}}
      
%-------------------------- Portada -----------------------%
\newcommand{\oldtitle}[1]{\def\OldTitle{#1}}
\newcommand{\changerationale}[1]{\def\ChangeWhy{#1}}

\newcommand{\changebox}{%
\begin{tcolorbox}[title={Historial de cambios del tema}]
\textbf{Título anterior:} \OldTitle\par
\textbf{Motivo del cambio:} \ChangeWhy
\end{tcolorbox}
}

%-------------------------- Author Font -----------------------%
    \newcommand{\authorfont}[1]{{\fontfamily{EBGaramond-TLF}\selectfont #1}}

%-------------------------- Anexos -----------------------%
    \makeatletter
    \newcounter{anexo}
    \renewcommand{\theanexo}{\Roman{anexo}}
    \newcommand{\anexo}[1]{%
      \clearpage
      \phantomsection
      \refstepcounter{anexo}%
      \noindent\textbf{Anexo \theanexo.- }#1\par\medskip
      \addcontentsline{stc}{section}{Anexo \theanexo.- #1}%
    }
    \makeatother

    %-------------------------- Lista de figuras (personal) -----------------------%
\makeatletter
\newcommand{\MyListOfFigures}{%
  \clearpage
  \phantomsection
  {\centering\bfseries\Large Índice de figuras\par}%
  \medskip
  % Entrada en nuestro índice de contenido (stc)
  \addcontentsline{stc}{section}{Índice de figuras}%
  % Recuperación del listado (sin cabecera propia)
  \begingroup
    \parindent=0pt\parskip=2pt
    \@starttoc{lof}%
  \endgroup
}
\makeatother

%-------------------------- Bibliografía -----------------------%
\makeatletter
% Contador para las referencias
\newcounter{myref}
% Cabecera de la bibliografía, entra en el índice 'stc'
\newcommand{\MyBibliography}{%
  \clearpage
  \phantomsection
  {\centering\bfseries\Large Bibliografía y referencias\par}%
  \medskip
  \addcontentsline{stc}{section}{Bibliografía y referencias}%
  \setcounter{myref}{0}%
}
\newcommand{\MyBibItem}[4]{%
  \refstepcounter{myref}%
  \noindent[\themyref]\ #1.\ \emph{#2}. #3. #4\par
}
\makeatother

%--------- Establecer cuatro niveles de índice --------%
    \setcounter{tocdepth}{4}   
    \setcounter{secnumdepth}{4}% numera hasta \paragraph

%%%%%%%%%%%%%%%%%%%%%%%%%%%%%%%%

\makeatletter

    %--------- Interlineado Global --------%
    \edef\GlobalBaseLineStretch{\baselinestretch}
    
    %--------- Espaciado de seccinones --------%
    \renewcommand\section{\@startsection{section}
      {1}{\z@}
      {2.5ex \@plus 1ex \@minus .2ex} % espacio antes
      {1.5ex \@plus .2ex}             % espacio después
      {\normalfont\Large\bfseries}    % formato del título
    }
    \renewcommand\subsection{\@startsection{subsection}{2}{\z@}
      {3ex \@plus 0.8ex \@minus .2ex}
      {1.5ex \@plus .2ex}
      {\normalfont\large\bfseries}
    }
    \renewcommand\subsubsection{\@startsection{subsubsection}{3}{\z@}
      {3ex \@plus 0.5ex \@minus .2ex}
      {1ex \@plus .2ex}
      {\normalfont\normalsize\bfseries}
    }
    \renewcommand\paragraph{\@startsection{paragraph}{4}{\z@}%
    {-2ex \@plus -0.5ex \@minus -.2ex}% espacio antes
    {0.8ex \@plus .2ex}% espacio después
    {\normalfont\normalsize\itshape}}% ← cursiva en lugar de negrita

    %--------- Ecuaciones --------%   
    % Variables globales para almacenar títulos y notas
    \gdef\leq@current@section{}%
    \gdef\leq@current@subsection{}%
    \gdef\leq@last@printed@section{}%
    \gdef\leq@last@printed@subsection{}%
    
    % Comando para añadir nota (invisible en el cuerpo)
    \newcommand{\eqnote}[1]{%
      \addcontentsline{leq}{leqnote}{#1}%
    }
    
    % Comando para ecuaciones
    \newcommand{\eqblock}[2]{%
      % Verificar si necesitamos imprimir el título de sección
      \ifx\leq@current@section\leq@last@printed@section\else
        \addcontentsline{leq}{leqsection}{\leq@current@section}%
        \global\let\leq@last@printed@section\leq@current@section
        \gdef\leq@last@printed@subsection{}% Reset subsección al cambiar sección
      \fi
      % Verificar si necesitamos imprimir el título de subsección
      \ifx\leq@current@subsection\leq@last@printed@subsection\else
        \ifx\leq@current@subsection\@empty\else
          \addcontentsline{leq}{leqsubsection}{\leq@current@subsection}%
          \global\let\leq@last@printed@subsection\leq@current@subsection
        \fi
      \fi
      % Ecuación
      \refstepcounter{eqnum}%
      \begin{equation*}
        #1\tag{E.\theeqnum}
      \end{equation*}%
      \addcontentsline{leq}{leqt}{[E.\theeqnum] -- #2}%
      \addcontentsline{leq}{leqm}{#1}%
    }
    
    %%% ----- Índice de ecuaciones ----------%%%
    % Tamaño para []
    \newcommand{\leqIndexFont}{\normalsize}
    \newcommand{\leqTitleFont}{\tiny}
    \newcommand{\leqNoteFont}{\tiny}
    % Cabecera de sección 
    \newcommand*\l@leqsection[2]{%
      \addvspace{25pt}% Mayor separación antes de nueva sección
      \noindent\leqIndexFont
      \makebox[\linewidth]{\rule{.38\linewidth}{0.4pt}\ \ \textbf{  \MakeUppercase{#1}}\ \ \rule{.38\linewidth}{0.4pt}}%
      \par\vspace{-10pt}
    }
    % Cabecera de subsección 
    \newcommand*\l@leqsubsection[2]{%
      \par\addvspace{15pt}%
      \noindent\leqIndexFont #1
      \par\vspace{-7pt}
    }
    % Título de ecuación 
    \newcommand*\l@leqt[2]{%
    \par\addvspace{10pt}%
      \par\noindent\leqTitleFont #1\par
    }
    % Miniatura de ecuación
    \newcommand*\l@leqm[2]{%
      \vspace{0.3\baselineskip}%
      \par\scriptsize\noindent\hspace*{1.5em}\(#1\)\par%
    }
    % Nota de ecuación 
    \newcommand*\l@leqnote[2]{%
      \vspace{0.2\baselineskip}%
      \par\noindent\leqNoteFont\hspace*{1.5em}\textbf{Nota.--} #1\par\vspace{0.5\baselineskip}%
    }
    % Generación del índice
    \newcommand{\mylistofequations}{%
      \clearpage
      \phantomsection
      % Título visual de la lista (ajústalo a tu gusto):
      {\centering\bfseries\Large Índice de ecuaciones\par}
      % Entrada en nuestro índice 'stc':
      \addcontentsline{stc}{section}{Índice de ecuaciones}%
      % Llamada a tu lista real (usa el nombre exacto de tu macro):
      \@starttoc{leq}%
    }
    
    %--------- Captura de títulos de sección --------%
    \let\leq@original@section\section
    \let\leq@original@subsection\subsection
    % Flag para saber si estamos en una sección asterisco
    \newif\ifLeqStarSection
    \LeqStarSectionfalse
    
    % Redefinir \section CON CONTROL DE ASTERISCO
    \renewcommand{\section}{%
      \@ifstar{\leq@section@star}{\@dblarg\leq@section@nostar}%
    }
    \def\leq@section@star#1{%
      \global\LeqStarSectiontrue
      \gdef\leq@current@section{#1}%
      \gdef\leq@current@subsection{}%
      \leq@original@section*{#1}%
      \global\LeqStarSectionfalse
    }
    \def\leq@section@nostar[#1]#2{%
      \global\LeqStarSectionfalse
      \gdef\leq@current@section{#2}%
      \gdef\leq@current@subsection{}%
      \leq@original@section[#1]{#2}%
    }
    % Redefinir \subsection
    \renewcommand{\subsection}{%
      \@ifstar{\leq@subsection@star}{\@dblarg\leq@subsection@nostar}%
    }
    \def\leq@subsection@star#1{%
      \global\LeqStarSectiontrue
      \gdef\leq@current@subsection{#1}%
      \leq@original@subsection*{#1}%
      \global\LeqStarSectionfalse
    }
    \def\leq@subsection@nostar[#1]#2{%
      \global\LeqStarSectionfalse
      \gdef\leq@current@subsection{#2}%
      \leq@original@subsection[#1]{#2}%
    }

    % ----- Restauración de valores Globales de estilo ----------%
    % Flags
    \newif\ifInFootnote
    \InFootnotefalse
    \pretocmd{\@footnotetext}{\InFootnotetrue}{}{}
    \apptocmd{\@footnotetext}{\InFootnotefalse}{}{}
    % Profundidad de listas (soporta anidadas)
    \newcount\MyListDepth \MyListDepth=0
    \AtBeginEnvironment{la}{\advance\MyListDepth by 1}
    \AfterEndEnvironment{la}{\advance\MyListDepth by -1}
    \AtBeginEnvironment{lnum}{\advance\MyListDepth by 1}
    \AfterEndEnvironment{lnum}{\advance\MyListDepth by -1}
    \AtBeginEnvironment{itemize}{\advance\MyListDepth by 1}
    \AfterEndEnvironment{itemize}{\advance\MyListDepth by -1}
    \AtBeginEnvironment{enumerate}{\advance\MyListDepth by 1}
    \AfterEndEnvironment{enumerate}{\advance\MyListDepth by -1}
    % Restauración diferida: hazlo al INICIO del próximo párrafo real
    \newif\if@pendingRestore
    \@pendingRestorefalse
    \newcommand{\RequestRestore}{\global\@pendingRestoretrue}
    \everypar{%
      \if@pendingRestore
        \global\@pendingRestorefalse
        \ifInFootnote\else
          \ifnum\MyListDepth=0
            \setlength{\parskip}{\GlobalParskip}%
            \setstretch{\GlobalBaseLineStretch}%
          \fi
        \fi
      \fi
    }
    % Al final de bloques:
    \AfterEndEnvironment{equation}{\RequestRestore}
    \AfterEndEnvironment{equation*}{\RequestRestore}
    \AfterEndEnvironment{la}{\RequestRestore}
    \AfterEndEnvironment{lnum}{\RequestRestore}
    % Al final de macros:
    \apptocmd{\eqblock}{\RequestRestore}{}{}
    \apptocmd{\imagenfit}{\RequestRestore}{}{}
    
\makeatother

% ===================== ToC propio con 4 niveles (único bloque) =====================
\makeatletter

% --- Escritores: añadimos una línea al .stc en cada cabecera numerada ---
% Guardamos las versiones actuales para llamar al comportamiento normal
\let\MyTOC@orig@section\section
\let\MyTOC@orig@subsection\subsection
\let\MyTOC@orig@subsubsection\subsubsection
\let\MyTOC@orig@paragraph\paragraph

% SECTION
\renewcommand{\section}{%
  \@ifstar{\MyTOC@section@star}{\@dblarg\MyTOC@section@nostar}%
}
\def\MyTOC@section@star#1{%
  \MyTOC@orig@section*{#1}% sin entrada al .stc
}
\def\MyTOC@section@nostar[#1]#2{%
  \MyTOC@orig@section[{#1}]{#2}% comportamiento normal (escribe al .toc)
  % Escribimos también nuestra línea al .stc (texto con \numberline)
  \addcontentsline{stc}{section}{\protect\numberline{\thesection}#2}%
}

% SUBSECTION
\renewcommand{\subsection}{%
  \@ifstar{\MyTOC@subsection@star}{\@dblarg\MyTOC@subsection@nostar}%
}
\def\MyTOC@subsection@star#1{%
  \MyTOC@orig@subsection*{#1}%
}
\def\MyTOC@subsection@nostar[#1]#2{%
  \MyTOC@orig@subsection[{#1}]{#2}%
  \addcontentsline{stc}{subsection}{\protect\numberline{\thesubsection}#2}%
}

% SUBSUBSECTION
\renewcommand{\subsubsection}{%
  \@ifstar{\MyTOC@subsubsection@star}{\@dblarg\MyTOC@subsubsection@nostar}%
}
\def\MyTOC@subsubsection@star#1{%
  \MyTOC@orig@subsubsection*{#1}%
}
\def\MyTOC@subsubsection@nostar[#1]#2{%
  \MyTOC@orig@subsubsection[{#1}]{#2}%
  \addcontentsline{stc}{subsubsection}{\protect\numberline{\thesubsubsection}#2}%
}

% PARAGRAPH
\renewcommand{\paragraph}{%
  \@ifstar{\MyTOC@paragraph@star}{\@dblarg\MyTOC@paragraph@nostar}%
}
\def\MyTOC@paragraph@star#1{%
  \MyTOC@orig@paragraph*{#1}%
}
\def\MyTOC@paragraph@nostar[#1]#2{%
  \MyTOC@orig@paragraph[{#1}]{#2}%
  % Si no numeras los \paragraph, cambia \theparagraph por texto plano sin \numberline
  \addcontentsline{stc}{paragraph}{\protect\numberline{\theparagraph}#2}%
}

% --- Impresor del índice propio (lee el .stc) ---
\newcommand{\MyTOC}{%
  \par\bigskip
  {\centering\bfseries\Large Índice de contenido\par}%
  \medskip
  \begingroup
    \parindent=0pt\parskip=2pt
    \@starttoc{stc}%
  \endgroup
}

\makeatother
% ===================== fin ToC propio =====================


%%%%%%%%%%%%%%


% --- Datos del tema ---
\title{Tema 3.A.18 -- Análisis de mercados (III)\\
\large Diferenciación de productos: la teoría de la competencia monopolística y otros desarrollos.}
\author{Marcelo Ortega}
\date{Febrero 2026}
\newcommand{\TemaCode}{3A12}

\oldtitle{Tema 3.A.17. La teoría de la competencia monopolística y la diferenciación de productos.}
\changerationale{
    El cambio de título busca destacar que el tema debería tratar los principales modelos explicativos de la diferenciación de productos; la teoría de la competencia monopolística y otros desarrollos fundamentales en la materia.
}

\usepackage{soul}\sethlcolor{yellow}% resaltado amarillo: \hl{texto} (Panel TCEE, ⌃H)
\begin{document}


\maketitle
\changebox

\newpage

% --- Página de resumen y modificaciones ---
\puntosclave{ %% Escribir puntos clave Apuntes CECO
    }
\vspace{1cm}

\modificaciones{
    \textbf{OJO!} En modelos de producto no se debería incluir BETRAND ni COURNOT con diferenciación (son más propios de la competencia oligopolística). Se puede añadir para completar DIXIT-STIGLITZ (como previos):

    SPENCE (1976): competencia monopolística con elección endógena de calidad (calidad como atributo del producto).
    
	LANCASTER (1979): enfoque de características del producto (variedades definidas en un espacio de atributos, no espacial).

    O el propio Modelo de CHAMBERLAN (antecesor de DIXIT-STIGLITZ)
    }

\newpage
\MyTOC
\newpage

\mylistofequations
\clearpage

\MyListOfFigures
\clearpage


\pagenumbering{arabic}
\setcounter{page}{1}


\section{Introducción}\label{intro}


\subsection{Contextualización}\label{context}


\subsubsection{Relevancia}\label{importance}


\subsubsection{Historia}\label{history}



\subsubsection{Problemática}\label{problem} 




\subsection{Estructura de la exposición}\label{structure}



\clearpage
\section{Competencia monopolística: rasgos definitorios}
\puntosclave{
    En el monopolio un producto, el equilibrio de mercado viene dado por la maximización del beneficio por parte del monopolista, siéndole indiferente para ello elegir precio o cantidad. Esta maximización se produce cuando iguala sus ingresos y costes marginales.
}
	
La teoría microeconómica neoclásica se basa en el empleo de bienes homogéneos para modelizar el comportamiento de los agentes económicos. Sin embargo, en la práctica se observa que muchas industrias, incluyendo la mayoría de las industrias de bienes de consumo, producen un número elevado de bienes similares, aunque diferenciados: lo que conocemos como variedades.\footnote{%
    Más exactamente, se entiende por bienes diferenciados aquéllos que son considerados sustitutivos cercanos entre sí desde el punto de vista del consumo (i.e., son sustitutivos imperfectos entre sí en la función de utilidad de los consumidores), lo cual depende en última instancia de los gustos o preferencias de los consumidores. Por contraposición, los bienes homogéneos se podrían definir como aquellos que son sustitutivos perfectos en las funciones de utilidad de los consumidores; esto es, la elasticidad cruzada de la demanda es infinita a igualdad de precios.
}

\subsection{Competencia monopolística: Diferenciación}
Los bienes pueden estar diferenciados en muchas características o dimensiones tales como su localización geográfica, la calidad, el diseño, la presentación, grado de servicios proporcionado con el producto (asistencia técnica, servicio postventa, garantía etc.).

Por lo tanto, con el término diferenciación se alude a dichas variaciones dentro de una clase de productos que los consumidores (o algunos de ellos) consideran como sustitutos imperfectos. La literatura económica distingue dos tipos básicos de diferenciación de productos:
\begin{la}
    \item \textbf{Diferenciación horizontal}. Según la cual las diferencias entre productos se derivan de la diversidad de preferencias individuales sobre las características de un bien (diseño, cercanía, etc.);\footnote{%
        Un ejemplo de diferenciación horizontal sería la compra de un jersey: unos consumidores preferirán un color y otros uno distinto. En este caso, las empresas ofrecerán distintas combinaciones de características dirigidas a diferentes tipos de consumidores. Esta diferenciación persigue segmentar a los clientes y ganar poder de mercado sobre segmentos de clientes específicos a través de la creación de fidelidad la marca.
    } y,
    \item \textbf{Diferenciación vertical}. Conforme a la cual las diferencias entre productos no se derivan de las preferencias individuales sobre una de las características de un bien, ya que todos los consumidores la valoran de igual modo (calidad, seguridad, fiabilidad), sino de la distinta disponibilidad a pagar por una cantidad de esa característica presente en el bien.\footnote{%
        Un ejemplo de diferenciación vertical sería la compra de un televisor: en general, todo consumidor prefiere una mejor calidad de imagen y sonido, pero cada uno tiene distinta disponibilidad a pagar por un televisor con una determinada calidad de imagen y sonido. Por lo tanto, ello presupone la existencia de criterios universalmente aceptados para evaluar la calidad.
    }  
\end{la}

A tenor de estas definiciones, la diferenciación horizontal se asocia a la variedad, mientras que la diferenciación vertical a la calidad; pudiendo en cualquier caso, considerarse la presencia de ambos tipos en cualquier entorno económico.\footnote{%
    Obsérvese que en la práctica no es infrecuente encontrar presentes en un producto a ambos tipos de diferenciación. Por ejemplo, cuando un consumidor elige el tipo de coche que puede permitirse, nos hallamos ante un problema de diferenciación vertical, pero una vez se ha decidido, la selección del modelo es básicamente una cuestión de diferenciación horizontal.
} 

\subsection{Competencia monopolística: variación conjentural}
\textcolor{red}{\textbf{OJO! AQUÍ ES IMPRESCINDIBLE METER LA PARTE DE INTERDEPENDENCIA ESTRATÉGICA Y VARIACIÓN CONJENTURAL NULA!!!}} ES LA APORTACIÓN MÁS IMPORTANTE DE LA TEORÍA (en términos comparativos)

El objeto de nuestra exposición consistirá en estudiar la diferenciación de productos, que podemos hallar en dos estructuras de mercado (de competencia imperfecta): la competencia monopolística y el oligopolio:\footnote{%
    Para ello, nos apoyaremos en los resultados que proporciona la teoría de la organización industrial, que siguiendo a TIROLE (1988) se puede definir como aquella rama de la teoría microeconómica que tiene por objeto el estudio el funcionamiento de los mercados e industrias, especialmente la estructura de la empresa (i.e., su organización interna) y el modo en que compiten unas empresas con otras. Precisamente este es el motivo principal por el cual se considera como una rama diferenciada: el énfasis que pone en el estudio de las estrategias de las empresas, rasgo que caracteriza la interacción en el mercado (competencia en precios, posicionamiento de producto, publicidad, investigación y desarrollo, etc.)
}
\begin{la}
    \item La competencia monopolística se define como aquella estructura de mercado que presenta las mismas características que la competencia perfecta, salvo que el bien que produce cada empresa no es homogéneo, sino diferenciado en alguna característica;\footnote{%
        Ejemplos de la realidad que cabría asimilar a la competencia monopolística son la oferta de restaurantes en una ciudad o de libros de un determinado género.
    } y,
    \item Por otro lado, el oligopolio se define como aquella estructura de mercado que contiene un número reducido de empresas. En concreto, se considerará el caso en que producen un bien diferenciado. 
\end{la}

\textcolor{red}{Dado el número de agentes en esta estructura de mercado: la interacción estratégica juega un papel central (esto es, cada empresa es consciente que el resultado de su decisión depende de las decisiones tomadas por sus competidores), que hace indispensable el empleo de la teoría de juegos (en su versión no cooperativa) para su análisis; y, cada empresa tiene cierto poder de mercado (parcial). Esto es, la capacidad para fijar precios supracompetitivos por sí sola, lo cual determina la aparición de beneficios extraordinarios.}

\subsection{Competencia monopolística: caracterización}
Compartirán, no obstante, las siguientes hipótesis de partida:
\begin{lnum}
    \item Agentes racionales. Los consumidores maximizan su utilidad y las empresas maximizan sus beneficios;
    \item \textcolor{red}{Información perfecta. No existen problemas de información (no es ni incompleta ni asimétrica),  es decir: los consumidores conocen las características del producto y as empresas conocen la demanda del producto; y,}
    \item Producto diferenciado. A diferencia de la situación de competencia perfecta, los bienes producidos por las empresas no son sustitutos perfectos en las funciones de utilidad de los consumidores (i.e. en las funciones de utilidad de los consumidores).
\end{lnum}

\paragraph{Competencia monopolística: variación conjentural nula}
En esta exposición, caracterizaremos la competencia monopolística en base a los siguientes supuestos:
	Muchos oferentes (potenciales) y muchos demandantes.
	Libre entrada en el largo plazo. Si a corto plazo existen beneficios extraordinarios, entran empresas hasta que se eliminan estos beneficios.
De este modo, a diferencia de la competencia perfecta [ver Tema 3.A.16] las empresas tendrán cierto grado de poder de mercado al ser monopolistas de su propia variedad (i.e. se enfrentarán a una curva de demanda que no será totalmente elástica), y por lo tanto, podrán fijar precios supracompetitivos por sí mismas (i.e. serán precio-decisoras y no precio-aceptantes sin que, en este caso exista equivalencia entre fijación de precios y cantidades, como sí había en el mercado de monopolio). 

\paragraph{Competencia olipolística: interdependencia estratégica}
Caracterizaremos esta estructura de mercado en base a los siguientes supuestos:
	Pocos oferentes y muchos demandantes. Al haber muchos demandantes se comportarán de forma precio-aceptante. Sin embargo, la existencia de un reducido número de oferentes les dotará de poder de mercado y ocasionará interdependencia estratégica. 
	Barreras de entrada. Esto permitirá la existencia de beneficios a largo plazo (a priori) y pueden deberse a: (i) comportamientos estratégicos por parte del monopolista que impiden la entrada a potenciales competidores; (b) restricciones legales o administrativas que conceden una situación de privilegio monopolista a una determinada empresa (p.e. concesión de licencias, derechos de patentes, imposición de barreras comerciales que excluyan competidores extranjeros); o, (iii) la existencia de una tecnología de producción que propicie la existencia de un número reducido de empresas, determinando así una situación de oligopolio natural. 
    
Es decir, a diferencia de la estructura de mercado de competencia monopolística, existirán barreras de entrada que provocarán que el número de empresas sea limitado (trabajaremos con duopolio por simplicidad analítica) y se produzca poder de mercado en el corto plazo (con posibilidad de horizonte infinito). 
Además, el hecho de que existan un número reducido de empresas  introduce un fenómeno que es diferencial respecto a otras estructuras de mercado: la interdependencia estratégica (esto es, las acciones de otros decisores a una empresa de referencia tienen impacto sobre la función objetivo de una empresa dada), que conduce al uso intensivo de la Teoría de Juegos en el estudio de la competencia oligopolística.  


2.3. Una aproximación general a la modelización de la diferenciación de productos
Podemos describir de manera general un modelo completo de diferenciación de productos, especificando: (i) el conjunto de posibles productos; (ii) la tecnología asociada a cada producto; (iii) los gustos o preferencias de los consumidores sobre el conjunto de posibles productos; y, (iv) el concepto de equilibrio de mercado empleado. 
Puesto que el estudio de un modelo de diferenciación de productos con un grado significativo de generalidad sería inmanejable, gran parte de la literatura incurre en fuertes supuestos simplificadores de un tipo u otro. A continuación, nos detendremos en los rasgos más sobresalientes de esta modelización.
2.3.1. La tecnología
En la práctica, se observa que el conjunto de bienes fabricados por las empresas pertenecientes a una misma industria constituye un subconjunto reducido del conjunto de posibles productos. Para explicar esta observación empírica, 
	Se ha rechazado que se deba al hecho de que los consumidores sólo demanden productos en dicho subconjunto;  y,
	Se ha acudido mayoritariamente a la ausencia de convexidad en los conjuntos de producción de los distintos bienes, lo cual se derivaría de: 
(i) la existencia o no de costes de desarrollo (por ejemplo, de l+D) y comercialización del producto; y, (ii) la  indivisibilidades en el capital fijo; 
Que implican costes medios decrecientes a lo largo de un rango inicial de la producción. 
La no convexidad en los conjuntos de producción se refleja normalmente mediante el supuesto de una función de costes sencilla con un coste fijo de entrada, y (por comodidad analítica) un coste marginal de producción constante que puede estar sujeto o no a una restricción de capacidad. 
2.3.2. Preferencias de los consumidores
Cabe distinguir dos enfoques básicos para modelizar la demanda de bienes diferenciados (esto es, las preferencias del consumidor):
2.3.2.1. Enfoque basado en direcciones: adress approach
Esta rama de la literatura postula una distribución de las preferencias de los consumidores a lo largo de un espacio continuo de parámetros que describen la naturaleza de los productos: distintos consumidores tendrán distintas localizaciones preferidas en este espacio y, por lo tanto, puede concebirse que los consumidores tienen distintas direcciones en este espacio.  En palabras, un bien debe considerarse una cesta de características y las características, a su vez, son los argumentos de la función de utilidad.  
Por lo tanto, la demanda agregada de bienes diferenciados viene dada por las distintas preferencias de los consumidores (o grupo de consumidores) con respecto a los parámetros que describen la naturaleza de los bienes.
2.3.2.2. Enfoque no basado en direcciones: non-address approach
La otra rama no se aparta de la teoría microeconómica neoclásica en la medida en que las preferencias de los consumidores vienen definidas sobre los propios bienes diferenciados (i.e., 'los bienes son bienes' y como tales entran en la función de utilidad de los consumidores). Dentro de esta rama, existen dos posibles formas de reflejar la compra de múltiples bienes diferenciados: 
	La primera supone que las preferencias agregadas por los bienes diferenciados se pueden reflejar mediante la ficción de un consumidor representativo, lo cual implica que todos los consumidores se consideran idénticos.  Este consumidor representativo valora la variedad en sí misma y prefiere consumir el mayor número de variedades de un mismo producto. Por lo tanto, detrás de la demanda de bienes diferenciados sólo hay una preferencia (deseo) por la variedad.  Los autores que introdujeron este enfoque fueron SPENCE (1976), y DIXIT y STIGLITZ (1977). 
	La segunda aproximación seguida supone diferencias en las preferencias individuales (con la peculiaridad de que el consumidor puede comprar más de una variedad del bien diferenciado) y deduce el comportamiento agregado a partir de motivos individuales (HART, 1985; PERLOFF y SALOP, 1985).  
Indpendientemente de la rama, la mayoría de los modelos han empleado alguna variante del supuesto de simetría de CHAMBERLIN, el cual en términos generales establece que todos los productos son sustitutivos entre sí en la misma medida. 
Tendremos oportunidad de abordar con mayor detenimiento ambos enfoques al hilo del estudio de las distintas estructuras de mercado en las que se encuentra presente la diferenciación de productos: el enfoque de consumidor representativo, en la moderna teoría de la competencia monopolística, y el enfoque de direcciones, en las teorías del oligopolio. 
NOTA.- En lo introducido reside la esencia de nuestra exposición:
	Los modelos de producto consideran que los consumidores tienen preferencias heterogeneas sobre las diferentes variedades del bien, pero éstas no condicionan de manera tan determinante su elección de consumir una u otra variedad de forma que podemos considerar que el gasto que éste dedicará a su consumo se distribuye de manera uniforme entre todas las variedades por lo que la interacción entre Empresas girará entorno al precio; que será, en última instancia, el factor que valorará el consumidor (puesto que la variedad tampoco le es tan crucial)  
	Los modelos espaciales consideran que los consumidores tienen diferentes disponibilidades a pagar atendiendo a la variedad que vayan a consumir. Es decir, el precio no constituye un factor tan determinante, mientras que sí lo es la diferente localización que pueda alcanzar el productor en el espacio multidimensional de características del producto, por ello es un factor clave también de este enfoque la cantidad de consumo del bien: se trata de un consumo discreto (compro una variedad, y ese es mi punto de saciedad). Su capacidad de posicionarse (y mantenerse) con su variedad en un espacio distinto al resto será el determinante clave que condicione su poder de mercado y, en consecuencia, su capacidad de extraer más o menos excedente del consumidor. 
2.3.3. El concepto de equilibrio
En cuanto al concepto de equilibrio, habría que distinguir entre el equilibrio a corto plazo, pudiéndose suponer que la entrada en la industria está cerrada y, en consecuencia, que el número de empresas es fijo, y el equilibrio a largo plazo, en el cual la entrada en la industria es una posibilidad real. 
En este mismo contexto, se deben especificar cuáles son las relaciones supuestas entre las empresas, siendo muy distinto que las empresas compitan en precios o en cantidades. Con independencia de la variable estratégica empleada, se supondrá que esta competencia no se repite en el tiempo.
Por todo lo que hemos expuesto hasta ahora, la decisión de optar por el uso de uno u otro tipo de diferenciación es en realidad fruto de la investigación previa que debe realizar el investigador y, más en concreto, sobre las preferencias de consumo y la existencia de barreras de entrada en el mercado que se quiera analizar. 

\clearpage
\section{Modelos de producto}
\puntosclave{
    La estructura de mercado de competencia monopolística se estudia básicamente a partir de la reformulación moderna del modelo pionero de CHAMBERLlN, debida a DIXIT y STlGLITZ.
    
	Como supuestos principales del modelo de DIXlT y STIGLITZ, cabe destacar la simetría (de sustitución), el empleo del enfoque de consumidor representativo, la consideración de dos sectores en la economía (el de bienes diferenciados y el resto de la economía) y la libertad de entrada.
	
    La diferenciación de productos resultante bajo este modelo se aj ustaría a la definición de diferenciación horizontal.
    
	El modelo de DIXIT y STlGLITZ aborda, en primer lugar, el problema de decisión del consumidor representativo, empleando una función de utilidad con dos argumentos:
	\begin{la}
	    \item Una función de subutilidad de elasticidad de sustitución constante, para el grupo de bienes diferenciados; y 
        \item Un bien compuesto o numerario, para el resto de bienes de la economía.
	\end{la}

    Para la resolución del problema, el consumidor elige primero entre el consumo de numerario y el consumo del grupo de bienes diferenciados y posteriormente reparte su consumo dentro del grupo de bienes diferenciados. Ello nos permitirá obtener las dos curvas de demanda de la teoría de la competencia monopolística: la demanda esperada y la demanda real de cada bien (y sus elasticidades).
	
    El equilibrio de mercado en el grupo de bienes diferenciados se obtiene empleando los dos conceptos de demanda y las condiciones de costes. Las soluciones explícitas requieren una especificación precisa de la curva de costes y de la función de utilidad
    
	El modelo de DIXIT y STIGLITZ no está exento de críticas, pero su influencia ha sido notable (comercio internacional, macroeconomía).
}

\textbf{INTRODUCIR AQUÍ UN EJEMPLO DE LAS MARCAS DE AGUA}

El enfoque de producto (Chamberlin, Bertrand y Cournot con diferenciación, y formalmente Dixit–Stiglitz) concibe la competencia monopolística como un equilibrio en el que cada empresa controla una variedad única, pero compite con muchas otras a través de la sustitución en preferencias. La fuente del poder de mercado no es la escasez de empresas, sino la imperfección en la sustitución entre productos. En Dixit–Stiglitz esto se formaliza mediante una función CES, donde la elasticidad de sustitución \sigma resume toda la estructura competitiva: cuanto mayor es \sigma, más intensa es la competencia y menor el margen. El objetivo teórico central de este enfoque es cerrar analíticamente el equilibrio monopolístico con entrada libre, mostrando cómo coexisten mark-ups positivos, variedad endógena y beneficios nulos en equilibrio. La diferenciación es abstracta y simétrica: no importa qué distingue a los productos, sino cuánto se sustituyen.

\subsection{Número de empresas fijo: diferenciación en COURNOT y BERTRAND}
A continuación trataremos las implicaciones que tiene la Teoría de la Competencia Monopolística sobre los tradicionales Modelos de COURNOT y BERTRAND, de competencia en un mercado de oligopolio con un número de empresas dado. 

En ambos modelos, la diferenciación de producto rompe la homogeneidad de bienes y permite que las empresas tengan poder de mercado, lo que implica precios por encima del costo marginal, pero la forma en que aparece ese poder de mercado difiere:
\begin{la}
    \item En BERTRAND con diferenciación, el poder de mercado surge por la elasticidad cruzada de la demanda. Es decir, si los bienes no son sustitutos perfectos, las empresas pueden fijar precios distintos del costo marginal; y,
    \item En COURNOT con diferenciación, el poder de mercado surge por el grado de sustitución en la demanda agregada. Es decir, la cantidad elegida por una empresa influye en el precio percibido de su variedad y en el de las demás.
\end{la}

Formalmente, se puede demostrar que:
p^B<p^C,q^B>q^C,π^B<π^C  

Sin embargo, la diferenciación (sustitución imperfecta) hace que cada empresa enfrente una demanda (precio) menos elástica que en el caso de bienes homogéneos, lo que permite márgenes sobre el coste marginal. 
Bajo demanda lineal y simetría, ambos modelos pueden generar resultados análogos en cuanto al equilibrio (precio, cantidad y bienestar), bien sea: (i) ajustando la elasticidad de sustitución (cuánto más diferenciados estén los productos menor grado de sustituibilidad habrá y mayor poder de mercado tendrá el competidor sobre el micro-monopolio de su variedad); o, (ii) ajustando el número de competidores (a mayor número de competidores, mayor producción agregada y, en consecuencia, menor precio).
En resumen:
	La contribución de DIXIT y STIGLITZ (1977) puede verse como una formalización continua y suavizada del Equilibrio Bertrand con productos diferenciados donde cada empresa fija su precio óptimo dado su demanda individual, pero la entrada libre determina el número de variedades.
	El modelo de Cournot, con diferenciación horizontal o vertical (véase SHUBIK y LEVITAN, 1980), ofrece un resultado estructuralmente análogo, donde el número de empresas y la producción agregada se determinan por la interacción entre rendimientos crecientes, costos fijos y elasticidad de sustitución.

3.2.1. Diferenciación horizontal en el Modelo de BERTRAND
Una de las variantes del modelo de ciudad lineal de HOTELLING, que trataremos más en profundidad en el siguiente apartado relativo a los Modelos Espaciales, proporciona una salida muy adecuada a la Paradoja propuesta por BERTRAND (1883):
3.2.1.1. Paradoja de BERTRAND
La Paradoja de BERTRAND describe una situación en la que dos jugadores (empresas) llegan a un estado de equilibrio de Nash donde ambas empresas cobran un precio igual al costo marginal, pese a que son un duopolio.  
En efecto, como estamos considerando un producto homogéneo, si una empresa fijara un precio más bajo se haría con toda la demanda. La otra empresa reaccionaría también bajando el precio. Este proceso seguiría hasta que el precio fuese igual al coste marginal, pues una vez en ese punto la empresa no tendría incentivos ni a bajarlo (obtendría pérdidas) ni a subirlo (perdería toda su cuota de mercado). Por tanto, el equilibrio en el modelo de BERTRAND es equivalente al de competencia perfecta (Paradoja de Bertrand). Esto sucede debido a los supuestos del modelo:
	Costes marginales iguales
	Producto homogéneo
	No hay limitaciones de capacidad
La paradoja de Bertrand aparece raramente en la práctica, ya que:
	Los productos reales casi siempre se diferencian de alguna manera aparte del precio (nombre comercial, si no en otra cosa), 
	Las empresas tienen limitaciones en su capacidad para fabricar y distribuir, y, 
	Dos empresas rara vez tienen costos idénticos.
Veamos como la diferenciación de productos proporciona una solución a esto

3.2.2.2. Modelo de HOTELLING con localización exógena: Decisión de precios
Considerando un continuo de consumidores que se distribuyen de forma uniforme en un intervalo [0,1].  Existen diferentes variedades pero únicamente relativas a la localización, esto es, consideramos el producto homogéneo y una situación de duopolio (existen dos empresas A y B) que da lugar a una competencia en precios (a la BERTRAND).
Cada consumidor deriva su utilidad de acuerdo a la siguiente función:
 Y, además, cada consumidor va a consumir una unidad de una variedad (demanda unitaria perfectamente inelástica).  
Estos supuestos se aproximan enormemente a los del Modelo de Bertrand ya que se considera que el producto es homogéneo, la función de costes es idéntica (C(q_{i}=c·q_{i})) y lo único que diferencia ambos productos es el precio: que ahora estará compuesto por una cantidad fija (por venir dado exógenemante, correspondiente al valor del bien) y una cantidad definida por los costes de transporte que le supone al consumidor (endógena).
Sea la distribución uniforme de los consumidores tal que:
 
Desde esta situación, podemos obtener las funciones de demanda de cada bien, hallando al consumidor indiferente (x̃):
 
Conociendo las demandas, podemos resolver el problema de maximización de beneficios de las empresas, suponiendo que compiten en precios y resolviendo, a través de las condiciones de primer orden, obtendremos las funciones de reacción de las empresas:
  
Finalmente, resolviendo el sistema de ecuaciones resultantes, obtendremos el equilibrio de Nash:
  
Nótese como el precio es igual al coste marginal c más un mark-up, t, que viene dado por los costes de transporte.
Si no hubiera costes de transporte (los bienes no estuvieran diferenciados), obtendríamos el resultado de BERTRAND (1883) [ver tema 3.A.19]. Por tanto, este modelo muestra como, manteniendo (prácticamente) los mismos supuestos, la diferenciación de productos supone una salida a dicha Paradoja. De modo que, pese a competir en precios, las empresas fijan precios supracompetitivos y obtienen beneficios positivos.

\subsection{Número de empresas endógeno: Modelo DIXIT-STIGLITZ}
Pero, ¿qué pasaría si el número de empresas fuese endógeno? Imaginémos que no existen barreras de entrada al mercado y, en consecuencia, cualquier productor que se lo proponga pudiese entrar al mercado ofreciendo una variedad de producto única y convirtiéndose en el monopolista de dicha variedad. ¿Cómo podríamos modelizar dicha situación?

Con una reformulación de DIXIT y STIGLITZ (1977), de la pionera teoría de CHAMBERLIN nos permite estudiar dicha situación. Partimos de los siguientes supuestos: (lo introducimos en tabla comparativa en los apuntes únicamente con el fin de ver las diferentes con el modelo de CAMBERLAIN)

\begin{center}
\setlength{\tabcolsep}{6pt}
\renewcommand{\arraystretch}{1.2}
\begin{tabular}{|p{0.30\linewidth}|p{0.30\linewidth}|p{0.30\linewidth}|}
\hline
\multicolumn{3}{|c|}{\textbf{CHAMBERLAIN}}\\
\hline
\multicolumn{2}{|c|}{\textbf{Supuesto de simetría}} & \\ 
\hline
\multicolumn{2}{|p{0.60\linewidth}|}{%
Todos los productos diferenciados pertenecientes a una misma industria son sustitutivos cercanos entre sí, en la misma medida. La elasticidad cruzada de estos productos es alta pero nunca infinita. Como consecuencia, las empresas no tienen que preocuparse por la decisión de especificación del producto, sino solamente por la entrada: cualquiera que sea la variedad que decidan producir, existe un mercado que tiene las mismas propiedades que las del resto de los bienes en el grupo de bienes diferenciados. Además: (i) las curvas a las que se enfrenta cada empresa son idénticas, pues las preferencias están equitativamente distribuidas; y, (ii) las curvas de costes de cada empresa también son idénticas.%
} &
De este modo, cada empresa monopoliza una variedad del producto, aunque comparte el mercado con el resto de la industria. En consecuencia, la empresa no es precio aceptante (i.e. tiene cierto grado de poder de mercado) y la curva de demanda a la que se enfrenta la empresa tiene pendiente negativa.\\
\hline
\multicolumn{2}{|c|}{\textbf{Libertad de entrada en el mercado (sin barreras de entrada ni salida).}} & \\
\hline
\multicolumn{3}{|p{0.90\linewidth}|}{%
Todas las empresas tienen igual acceso a la tecnología de producción más eficiente. Hay un gran número de empresas en el mercado, lo que implica que: (i) el impacto de cada empresa sobre sus rivales es despreciable (i.e. ausencia de interacción estratégica); (ii) cada empresa o producto no tiene un bien ``cercano'' directo en el espacio de productos, dado que compite por igual con todos (supuesto de simetría).%
}\\
\hline
\multicolumn{2}{|c|}{\textbf{Movilidad perfecta de factores de producción}} & \\
\hline
\multicolumn{3}{|p{0.90\linewidth}|}{%
Existe competencia perfecta en el mercado de los factores productivos, de modo que éstos se remuneran según su productividad marginal.%
}\\
\hline
\multicolumn{2}{|c|}{\textbf{Economía de dos sectores}} & \\
\hline
\multicolumn{3}{|p{0.90\linewidth}|}{%
La economía consta de dos sectores:
\begin{itemize}\setlength\itemsep{0pt}
  \item[(a)] un sector de productos diferenciados, sustitutivos entre sí; y
  \item[(b)] el resto de la economía, que se agregaría en un único bien compuesto o numerario.
\end{itemize}%
}\\
\hline
\multicolumn{3}{|c|}{ADICIONALES: \textbf{DIXIT--STIGLITZ}}\\
\hline
\multicolumn{2}{|c|}{\textbf{Enfoque de consumidor representativo}} & \\
\hline
\multicolumn{3}{|p{0.90\linewidth}|}{%
Existe un único consumidor representativo (i.e., no existe heterogeneidad en las rentas y los gustos) y este consumidor demanda todas las variedades existentes de un bien diferenciado. Por lo tanto, detrás de la demanda de bienes diferenciados sólo hay una preferencia (deseo) por la variedad.%
}\\
\hline
\end{tabular}
\end{center}

La empresa tiene cierto grado de poder de mercado que se proyecta como la capacidad para fijar el precio de su producto con un \textit{mark-up}, pero sólo de modo parcial, dado que se enfrenta a la competencia de los sustitutivos cercanos ofrecidos por las otras empresas. 

\subsubsection{Consumidor: Preferencia por la variedad}
Consideremos que el consumidor representativo presenta una Función de Utilidad sobre dos bienes, el numerario y el no numerario, $V(\cdot)$. Suponiendo que la función de ubutilidad ($V(\cdot)$) es simétrica y de tipos CES,\footnote{%
    Recuérdese que la elasticidad de sustitución entre dos bienes recoge el porcentaje de variación en la proporción en la que se consume un par de bienes ante variaciones porcentuales de su precio relativo. La función de subutilidad CES permite reflejar el supuesto de simetría.
} entonces estaríamos ante la siguiente especificación:

\eqblock{
\begin{aligned}
    &U(m,V(x_1,x_2,\dots,x_I))\quad \text{donde,}\;V(\cdot)\equiv \text{CES y simétrica}\\[0.5em]
    &\Longrightarrow V(\cdot)\equiv\left(\sum_{i=1}^Ix_i^\rho\right)^{1/\rho}\qquad \text{con,}\;\rho\in(0,1)\\[2em]
    &\boxed{u=U\left(m,\left[\sum_{i=1}^Ix_i^\rho\right]^{1/\rho}\right)}
\end{aligned}
}{Especificación de la función de utilidad: CES y simetría (separable). Competencia monopolística (variedades endógenas): DIXIT-STIGLITZ}

Una vez conocemos la estructura de la Función de Utilidad, podemos presentar el problema de elección del consumidor representativo. El vector de consumo de equilibrio del consumidor representativo debe resolverse mediante el siguiente problema de maximización sujeto a la restricción de recursos o dotación inicial del bien numerario, $m_0$:

\eqblock{
\begin{aligned}
    &\boxed{
    \begin{aligned}
        &\max U=U\left(m,\;\left\{\sum_{i=1}^Ix_i^\rho\right\}^{\frac{1}{\rho}}\right)\\[1em]
        &\text{s.a.}\quad m+\sum_{i=1}^Ip_i\cdot x_i\leq m_0
    \end{aligned}
    }\\[1em]
    &\text{(Solución: 2 Etapas)}\\[0.5em]
    &\Longrightarrow
    \left\{\begin{aligned}
        &\underset{\text{\scriptsize $S(\cdot)\;\;$ es el porcentaje del gasto que se destina a no numerario}}{\text{(1) \% $m_o$ en no numerario:}}&&m=m_0\cdot [1-S(p)]\Rightarrow \;\boxed{g=m_0\frac{S(p)}{p}}\\[0.5em]
        &\text{(2) Cantidad de cada variedad:}&&\boxed{x_i=g\cdot\left(\frac{p}{p_i}\right)^{\sigma}}\qquad \text{donde, }\;\sigma=\frac{1}{1-\rho}
    \end{aligned}\right.
\end{aligned}
}{Problema de maximización del consumidor: Utilidad. Competencia monopolística (variedades endógenas): DIXIT-STIGLITZ}

Dado que $U(\cdot)$ es una función de utilidad separable, el problema de optimización del consumidor puede resolverse en dos pasos separados: (i) primero se determina la asignación óptima del ingreso para cada subgrupo;\footnote{%
    La función $S(\cdot)$ depende de: (i) la forma precisa que adopte la fucnión de utilidad $U(\cdot)$; y, (ii) el ínidice de precios $p$ de los bienes diferenciados.
} y, (ii) luego las cantidades dentro de cada subgrupo.\footnote{%
    El supuesto de la separabilidad de la función de utilidad provoca que el porcentaje de gasto al grupo de productos diferenciados y al numerario dependa únicamente del índice de precios del grupo de bienes diferenciados (puesto que el precio del bien numerario es la unidad).
}

\paragraph{Condición CES: Simplificación}
Ahora bien, como hemos impuesto la elasticidad constante de sustitución (preferencias simétricas y homotéticas) podemos simplifcar enormemente nuestros resultados. 

Considerando que existen $I$ variedades diferentes, por la estructura de la función de subutilidad, el consumidor no tiene ningún motivo para discriminar entre variedades, por lo que se generarán tres implicaciones relevantes en equilibrio: (i) el gasto se distribuirá uniformemente entre todas ellas (pues el consumidor prefiere consumir el máximo de variedades posibles); (ii) el precio será igual para todas ellas, puesto que una variación infinitesimal del precio haría al consuimdor renunciar a esa variedad por su encarecimiento relativo; y, (iii) la cantidad producida es igual para todas. Por tanto, podemos reformular las ecuaciones como:

\eqblock{
\begin{aligned}
    &\text{Como,}\;V(\cdot)\equiv
    \left\{\begin{aligned}
        &\underset{\text{\scriptsize Entre variedades simetría}}{\text{Preferencias simétricas}}\\
        &\underset{\text{\scriptsize Preferencias Homotéticas}}{\text{Subutilidad CES}}
    \end{aligned}\right.
    \underset{\text{\scriptsize En equilibrio}}{\Longrightarrow}\;
    \left\{\begin{aligned}
        &p_i=p,&&\forall i\\
        &x_i=x,&&\forall i
    \end{aligned}\right\\[1em]
    &\text{\scriptsize
    $\Rightarrow\;\left\{\begin{aligned}
        &\text{Reescribimos,}\\[0.5em]
        &V=\left[\sum_{i=1}^Ix_i^\rho\right]^{1/\rho}=(x^\rho I)^{1/\rho}=x\cdot I^{1/\rho}\\[0.5em] 
        &\sum_{i=1}^Ip_ix_i=\sum_{i=1}^{I}xp=xpI
    \end{aligned}\right\}$}
    \overset{\text{\tiny Resolviendo}}{\Longrightarrow}
    pIx=m_0S(p)\Longrightarrow \underset{\substack{\text{\scriptsize Total consumido de variedades}\\\text{\scriptsize como función de gasto total}\\\text{\scriptsize y número de variedades}}}{\boxed{x=\frac{m_0S(p)}{p\cdot I^{1/\rho}}}}
\end{aligned}
}{Problema de maximización del consumidor: Función tipo CES (preferencias simétricas y homotéticas). Competencia monopolística (variedades endógenas): DIXIT-STIGLITZ}
  
Donde el primer término del lado derecho es la elasticidad de la participación del grupo de productos diferenciados en el gasto total del consumidor con respecto al índice de precios. El hecho de que la utilidad aumente con la variedad puede constatarse recordando g (tal y como viene expresado en (3)) y analizando cómo p (tal y como viene expresado en (5)) evoluciona con n, de acuerdo con lo recogido en el gráfico 1 (suponiendo que $p*= 1$):
 
Con un gasto constante, el índice de precios p disminuye rápidamente, por lo que la utilidad aumenta. Este efecto es más pronunciado para p cercanos a O. En cambio, para p cercanos a 1 , todas las variedades son sustitutivas cercanas y, por lo tanto, la introducción de otra variedad similar solo aumenta la utilidad de forma reducida.


Conviene hacer algunas observaciones sobre la función de demanda y las preferencias CES.

\paragraph{Elasticidades-renta unitarias}
Lo primero que debemos introducir es la elasticidad renta de la demanda de variedades, que nos permite medir cómo responde el consumo de cada variedad a un aumento proporcional de la renta, manteniendo precios constantes. Esto es, nos indica si los bienes diferenciados son normales, inferiores o de lujo dentro del agregado.

\eqblock{
\begin{aligned}
    &\varepsilon_{x_i,m}=\frac{\partial x_i}{\partial m}\cdot\frac{m}{x_i} &&\underset{x_i=\frac{m_0S(p)}{pI}}{\overset{\substack{\text{\scriptsize Con la especificación}\\\text{\scriptsize del problema}}}{\Longrightarrow}}&&\underset{\text{\scriptsize $S(p)$ independiente de $m_0$}}{\boxed{\varepsilon_{x_i,m}=1}}
\end{aligned}
}{Elasticidad-demanda renta para cada variedad. Competencia monopolística (variedades endógenas): DIXIT-STIGLITZ}

Las preferencias CES implican homoteticidad: un aumento proporcional de la renta incrementa proporcionalmente el consumo de todas las variedades. Por tanto, bajo esta especificación, todas las variedades son bienes normales con elasticidad-renta unitaria. Cuando la renta aumenta un 1\%, el consumo de cada variedad aumenta exactamente un 1\%, manteniéndose inalterada la estructura del consumo. Esta propiedad es clave para Dixit–Stiglitz porque permite: (i) agregar consumidores sin distorsión; (ii) analizar entrada sin preocuparse por efectos redistributivos; y, (iii) separar limpiamente decisiones de renta y de sustitución.

\paragraph{Elasticidad-precio de la demanda de una variedad}
Por otro lado, tenemos la elasticidad demanda precio de una variedad concreta, que nos permite medir la sensibilidad del consumo de una variedad a su propio precio, manteniendo constantes el resto de precios y la renta. 

\eqblock{
\begin{aligned}
    &\varepsilon_{x_i,p_i}=\frac{\partial x_i}{\partial p_i}\cdot\frac{p_i}{x_i} &&\underset{x_i=g\cdot\left(\frac{p}{p_i}\right)^\sigma}{\overset{\substack{\text{\scriptsize Con la especificación}\\\text{\scriptsize del problema}}}{\Longrightarrow}}&&\underset{\text{\scriptsize Sustitución \textbf{entre variedades}}}{\boxed{\varepsilon_{x_i,p_i}=-\sigma=-\frac{1}{1-\rho}}}\\[2em]
\end{aligned}
}{Elasticidad-precio de la demanda de cada variedad (precio propio). Competencia monopolística (variedades endógenas): DIXIT-STIGLITZ}

Dadas las especificacioes del problema, cada variedad enfrenta una demanda de elasticidad constante. Como veremos, esto implica que el markup del monopolista es constante y, por tanto, que el entorno de demanda es compatible con precios fijados con mark-up sobre el coste marginal (como veremos más adelante).
 
\paragraph{Elasticidad de sustitución: Intersectorial e Intrasectorial}
Por último, tenemos:
\begin{la}
    \item \textbf{Elasticidad del agregado respecto al precio común}. La elasticidad del agregado de bienes diferenciados respecto a su precio común, igual que en el modelo original de CHAMBERLAN, nos permite saber qué sucede cuando el índice de precios aumentan simultáneamente y en la misma proporción, donde no hay sustitución entre variedades si no sustitución con el numerario. Esta elasticidad, por tanto, no viene fijada por $\sigma$ (que mide cambios ante variaciones de precios relativos) y no es una propiedad de la subutilidad CES, sino del problema de asignación de gasto ($S(p)$), es decir, del margen intersectorial: cuánto reduce el consumidor su gasto en el conjunto de bienes diferenciados cuando su precio común sube frente al numerario; y,
    \item \textbf{Elasticidad-precio intrasectorial o variedades}. La elasticidad de sustitución entre variedades, que nos permite captar qué facilidad tiene el consumidor para sustituir una variedad por otra cuando cambian los precios relativos. Mide el grado de diferenciación efectiva. Este es la verdadera diferencia respecto al modelo de CHAMBERLAIN, ya que formaliza explícitamente y con microfundamentos claros la sustitución entre variedades mediante una elasticidad constante ($\sigma$), separándola limpiamente de la sustitución entre el agregado de diferenciados y el numerario.\footnote{%
        CHAMBERLAIN intuía esta distinción, pero no podía cerrarla analíticamente ni definir una elasticidad bien comportada entre variedades.
    }
\end{la}

\eqblock{
\begin{aligned}
    &\text{Desde CHAMBERLAIN:}\\[1em]
    &\varepsilon_{x,p}=\frac{\partial \ln{x}}{\partial \ln{p}}\qquad {\Longrightarrow}\qquad \underset{\text{\scriptsize Sustitución \textbf{intersectorial}}}{\boxed{\varepsilon_{x,p}=\frac{\partial \ln{S(p)}}{\partial p}-1}}\\[2em]
    &\text{Novedad: DIXIT-STIGLITZ}\\[1em]
    &\sigma_{i,j}=\frac{\partial \ln{(x_i/x_j)}}{\partial \ln{(p_j/p_i)}}\cdot\frac{p_i}{x_i}\quad \overset{\substack{\text{\scriptsize Con la especificación}\\\text{\scriptsize del problema}}}{\Longrightarrow}\quad\underset{\text{\scriptsize Sustitución \textbf{intrasectorial}: variedades}}{\boxed{\sigma_{i,j}=\sigma=\frac{1}{1-\rho}}}\\[1em]
    &\Rightarrow\substack{\text{Grado de}\\\text{diferenciación}}
    \left\{\begin{aligned}
        &\sigma\to\infty&&\Longrightarrow\text{Sustitutivos perfectos}&&\Longrightarrow\text{Competencia intensa ($\sim$ Perfecta)}\\[0.5em]
        &\sigma\to1&&\Longrightarrow\text{No sustitutivos}&&\Longrightarrow\text{Poder de mercado ($\sim$ Monopolio)}
    \end{aligned}\right.
\end{aligned}
}{Elasticidad-precio de la demanda de cada variedad (precio propio). Competencia monopolística (variedades endógenas): DIXIT-STIGLITZ}

La elasticidad intersectorial o elasticidad de sustitución intersectorial, explica por qué la demanda real de una variedad es, bajo simetría, simplemente una fracción fija de la demanda agregada del sector: cada empresa mantiene su cuota y lo que cambia es el tamaño del mercado. Por su parte, la elasticidad de sustitución entre variedades, nos permite saber cuánta presión competitiva viene asociada a la estructura de preferencias que conforman la demanda de variedades.

En consecuencia, el modelo DIXIT-STIGLITZ separa analíticamente dos márgenes distintos de ajuste de la demanda: (i) intersectorial, que determina el tamaño del mercado del sector frente al numerario; e, (ii) intrasectorial, gobernado por $\sigma$, que determina la intensidad competitiva entre variedades (variedad contra variedad) y el mark-up de equilibrio.

\paragraph{Índice de precios: Deflactor del bienestar}

\textcolor{red}{\textbf{OJO!}} -- Todo este apartado esta \textit{incorrecto} debe reformularse para comprenderse bien. Se puede encontrar el desarrollo concpetual en los anexos [\authorfont{Ver Anexo \ref{anx:deflactor_bienestar}}]

¿Cómo afecta el número de variedades al bienestar del consumidor? Este es un concepto tremendamente innovador de la modelización de la competencia monopolística a la DIXIT-STIGLITZ. Como vemos, el problema dual del consumidor nos proporciona un índice de precios que, conceptualmente, nos indica el coste mínimo necesario para adquirir una unidad real de utilidad en términos de la cesta de variedades. Como hemos visto, la utilidad en el modelo se adquiere mediante el agregado de consumo en variedades que el consumidor puede permitirse: si el gasto mínimo necesario para adquirir una unidad de $V(\cdot)$ es $q(\mathbf p)\equiv \mathbf p$, entonces, dado un gasto $E=m_0S(\mathbf p)$ la cantidad total del agregado que puede comprar será $V(\cdot)=\frac{E}{\mathbf p}$. Como sabemos que el índice de precios es una función del número de variedades y la elasticidad de sustitución entre variedades, podemos ver como este coste mínimo cambia en función de sus parámetros::

\eqblock{
\begin{aligned}
    &\text{Índice de precios:}&&q(\mathbf p)=\\[1em]
    &\underset{\substack{\text{\scriptsize Dada nuestra}\\\text{\scriptsize especificación}}}{\Longrightarrow}
    \text{\scriptsize$\left\{\begin{aligned}
        &p_i=p=1,\;\;\forall i\\
        &x_i=x,\;\;\forall i
    \end{aligned}\right\}$}
    &&q(\mathbf p)=\left(\sum_{i=1}^Ip^{\frac{\rho}{(\rho-1)}}\right)^{\frac{(\rho-1)}{\rho}}= p\cdot (I)^{\frac{(\rho-1)}{\rho}}\Longrightarrow\underset{\text{\scriptsize Índice de precios de variedades}}{\boxed{q(\mathbf p)=p\cdot (I)^{-\frac{1}{\sigma}}}}
\end{aligned}
}{Índice de precios: Intensidad competitiva. Competencia monopolística (variedades endógenas): DIXIT-STIGLITZ}

Bajo el supuesto de simetría (todas las variedades se venden al mismo precio) el índice de precios adopta una forma especialmente simple, que permite visualizar con claridad cómo el número de variedades y el parámetro $\rho$ (o, equivalentemente, $\sigma$) afectan al precio efectivo del sector:

\imagenfit{IndicePrecios_DIXIT_STIGLITZ.png}{Índice de precios $q(\mathbf p)$ como función del número de variedades $I$ (asumiendo precio constante, $p_i=p=1$)}{Apuntes ICEX-CECO. Tema 3.A.18}

Como vemos, bajo este marco añadir variedades no reduce precios individuales, sino que reduce el precio efectivo del agregado porque el consumidor puede redistribuir su gasto entre más opciones: la competencia relevante aquí es intrasectorial (entre variedades), no intersectorial. La pendiente y curvatura dependen crucialmente de $\rho$ (o, equivalentemente, de $\sigma=1/(1-\rho)$):
\begin{la}
    \item \textbf{Alto grado de sustituibilidad}. Cuando $\rho$ es bajo (alta $\sigma$, variedades muy sustitutivas), el índice $q$ cae muy rápido al aumentar $I$: cada nueva variedad aporta un fuerte beneficio competitivo, reduciendo de forma sustancial el precio efectivo del sector; y,
    \item \textbf{Bajo grado de sustituibilidad}. En cambio, cuando $\rho$ es alto (baja $\sigma$, variedades poco sustitutivas), el descenso de $q$ es mucho más lento: las nuevas variedades aportan diversidad, pero generan poca presión competitiva, de modo que el precio del sector apenas baja.
\end{la} 

La implicación clave es que el modelo separa nítidamente variedad de poder de mercado. Aumentar $\uparrow \#I$ siempre mejora el bienestar vía un menor índice de precios ($\downarrow q(\mathbf p)$), pero la magnitud del efecto depende de la elasticidad de sustitución. Esto explica por qué, en equilibrio de competencia monopolística, puede coexistir una gran variedad de productos con márgenes positivos: el tamaño del mercado (intersectorial) y la intensidad competitiva (intrasectorial, $\sigma$) son márgenes distintos.

\subsubsection{Productor: Competencia monopolística}
Ahora que conocemos la perspectiva de la demanda, pasemos a analizar el problema del productor.

Asumiendo que todas las empresas presentan la misma estructura de costes con rendimientos crecientes a escala ($F_j$, costes fijos) y que cada empresa produce una única variedad.\footnote{%
    La función de costes de la empresa tipo puede entenderse como: $CT(x_i)=C(x_i)+F$

Por tanto, no hay economías de alcanca. Recordar que todas tienen la misma estructura de costes por lo que no hay subíndices.
} Podemos formular el problema de maixmización del bemeficio del monopolista como:

\eqblock{
\begin{aligned}
    &\max\quad \Pi_j(x_i)=p(x_i)x_i-cx_i-F\\[1em]
    &\text{CPO}\;\Longrightarrow\;\underset{\text{\scriptsize Monopolista de su propia variedad}}{{IMg}_{x_i}={CMg}_{x_i}}\\[2em]
    &\Longrightarrow\;\underset{\text{\scriptsize Ingreso marginal}}{p\left[1-\frac{1}{\sigma}\right]}=c\;\Rightarrow\;p\left[1-\frac{1}{|\varepsilon_{x^d_i,p_i}|}\right]=c\\[1em]
    &\frac{p-c}{p}=\frac{1}{\sigma}\;\Longrightarrow\;
    \boxed{p^*=\frac{c}{\rho}=\frac{c}{(\sigma-1)/\sigma}=\underset{\substack{\text{\scriptsize Mark-up sobre}\\\text{\scriptsize coste marginal}}}{\mu\cdot c}}
\end{aligned}
}{Problema de maximización del productor: Beneficios. Competencia monopolística (variedades endógenas): DIXIT-STIGLITZ}
 
Como se puede apreciar, asumimos (por simplicidad) que cada empresa considera dado el índice agregado de precios. Es decir, no tienen en cuenta ni el comportamiento del resto de empresas ni la influencia que una misma pueda tener sobre el ínidice agregado de precios: la empresa considerará el precio estipulado por el índice agregado.\footnote{
Nótese que, de esta forma, aunque ligeramente diferente al escenario de mercado monopolístico, el resultado es el mismo: el monopolista (de su variedad) maximiza los beneficios ajustando vía cantidades, no vía precios. Es decir, producirá la cantidad de output que, dado el nivel de precios, maximice sus beneficios. 
} 

\paragraph{Condición de libre entrada}
Pero entonces, ¿qué información nos proporciona a título individual el problema anterior? A priori, no nos proporciona información más allá de la condición óptima. Para conocer las decisiones que debe tomar el monopolista debemos recurrir a otra condición muy importante del modelo: la libre entrada. Como sabemos, el modelo asume que no existen barreras a la entrada y, por tanto, cuando una empresa entra al mercado y comienza a producir una nueva variedad los consumidores desviarán parte del gasto que antes destinaban a las variedades existentes a comprar el nuevo bien. En consecuencia, la cantidad vendida de cada variedad disminuye, al igual que los beneficios a título individual.

De esta forma, la condición de libre entrada nos permite conocer la producción de cada variedad en equilibrio, puesto que el beneficio a título individual debe ser nulo:

\eqblock{
\begin{aligned}
    &\text{Condición de libre entrada}\Longrightarrow\;
    \left\{\begin{aligned}
        &\text{Si,}\;\;\Pi_j>0\Rightarrow &&\uparrow \#X=(x_1,\dots,x_I)\\[0.5em]
        &\text{Si,}\;\;\Pi_j<0\Rightarrow &&\downarrow \#X=(x_1,\dots,x_I)\\[0.5em]
        &\text{Si,}\;\;\Pi_j=0\Rightarrow &&\text{Eq:}\;\#X=(x_1,\dots,x_I)
    \end{aligned}\right.\\[1em]
    &\Pi_j(x_j)=\quad p^*x_j-cx_j-F=0\iff (p^*-c)x_j=F\Rightarrow\;x_j=\frac{F}{p^*-c}\\[1em]
    &\underset{\text{\scriptsize De una variedad determinada}}{\text{Cantidad producida:}}\qquad \boxed{x_j^*=\frac{F}{c}\cdot(\sigma-1)}
\end{aligned}
}{}
 
Por tanto, la cantidad óptima de variedad producida en equilibrio dependerá de:
\begin{la}
    \item \textbf{El cociente entre costes fijos y marginales} ($\frac{F}{c}$). Cuanto mayores sean los costes fijos en relación a los costes marginales más se podrán aprovechar las economías de escala. Se del determinante principal ya que es íntegramente parte de la decisión del productor; y,
    \item \textbf{Preferencia por la variedad} ($\sigma$). En función del grado de amor por la diversidad que presenten los consumidores, el equilibrio se alcanzará con menos o más cantidades producidas de producto. 
\end{la}

Esto identifica de manera única un equilibrio si el lado izquierdo es una función monótona de $I$, lo cual ocurre si la elasticidad respecto a n tiene signo determinado. Se asume que es negativa, ya que la cantidad consumida de cada variedad disminuye cuando hay más variedades disponibles. 

\subsubsection{Equilibrio de mercado}
El equilibrio en este modelo se determina mediante tres condiciones:
	Vaciamiento de mercados. En equilibrio, la cantidad demandada de cada variedad se iguala a la cantidad producida de ésta.
	Maximización de beneficios. Las empresas maximizan sus beneficios de acuerdo con su función de demanda individual (función de demanda de la variedad que producen, i)(13);
	Condición de libre entrada en el largo plazo. Como esto genera beneficios extraordinarios, existen incentivos para que nuevas empresas entren al mercad y, por lo tanto, las cantidades de la variedad i, se ajustan hasta que la empresa marginal simplemente cubre los gastos y no obtiene beneficios (beneficios nulos).

En el corto plazo, La empresa elige una cantidad donde el ingreso marginal sea igual al coste marginal y fija el precio de la variedad en función de la curva de demanda en ese punto [Imagen .a)]. En este punto, cabe la posibilidad de que se den beneficios o pérdidas extraordinarios. Por el contrario, en el largo plazo, debido a la libre entrada y salida de empresas, los beneficios positivos inducirían a la entrada de nuevas empresas (y las pérdidas a la salida de las ya instaladas) hasta que éstos desaparecieran. A medida que nuevas empresas entran en el mercado, el gasto total en la industria se distribuye entre una mayor cantidad de variedades y, en consecuencia, la curva de demanda de cada variedad se desplazará hacia la izquierda (market crowding effect) hasta que sea tangente a la curva de costes medios.  
	 
En cualquier caso, debe quedar claro, que en el modelo no hay dinámica de entrada o de salida, sino que se llega directamente al equilibrio del gráfico de la derecha [Imagen 7.c].




\clearpage
\section{Modelos espaciales}
\puntosclave{
    El estudio del monopolio multiproducto (esto es, un mismo monopolista, produce conjuntamente distintos bienes en régimen de monopolio) depende crucialmente de las relaciones que se establezcan entre las demandas y costes de producción de los bienes que se producen conjuntamente.
}

El enfoque de producto (Chamberlin, Bertrand y Cournot con diferenciación, y formalmente Dixit–Stiglitz) concibe la competencia monopolística como un equilibrio en el que cada empresa controla una variedad única, pero compite con muchas otras a través de la sustitución en preferencias. La diferenciación, por tanto, es abstracta y simétrica: \textbf{no importa qué distingue a los productos, sino cuánto se sustituyen.}

Por ejemplo, en el mercado de telefonía existen numerosos competidores que sin duda compiten en régimen de competencia monopolística. Ahora bien, las variedades ofrecidas no son precisamente sustitutivas desde el punto de vista del consumidor. Los consumidores persiguen diversos tipos de características que pueden ser (o no ofrecidos) por una o varias empresas. Si analizásemos la competencia en diferentes \textbf{espacios de características} podríamos ver que existen consumidores cuyas características están muy localizadas en el entorno cercano de un oferente de mercado (Apple, por ejemplo) mientras que existen otros grupos de consumidores que persiguen características cuyo lugar espacial esta ocupado por diversos oferentes (HUAWEI, MIU, etc.). Este tipo de mercados no pueden ser estudiados desde la perspectiva anteriormente presentada y, por ello, existe una vertiente teórica de la competencia monopolística que pretende dar respuesta a estos fenómenos.

El enfoque espacial parte precisamente de esta concepción: la competencia monopolística surge porque los consumidores tienen preferencias heterogéneas localizadas, y las empresas se diferencian al ocupar posiciones distintas en un espacio de características. El poder de mercado procede de la proximidad local y de los costes de desplazamiento. Aquí, la competencia no es global ni simétrica: cada empresa compite intensamente con sus vecinos cercanos y débilmente con las lejanas. 

\subsection{Número de empresas fijo: Modelo ciudad lineal}

\textbf{OJO! Aquí hay que dejar claro que se trata de un modelo más propio del ámbito de la competencia oligopolística que al de la competencia monopolística. Lo introduciemos aquí con el fin de servir como puente para el Modelo de Ciudad Circular que, en su modalidad de libre entrada sí constituye un modelo de competencia monopolística:

El modelo de ciudad circular de Salop es el puente entre ambos mundos. Para un número pequeño de empresas sigue siendo un oligopolio espacial. Sin embargo, cuando se introduce entrada libre y se analiza un equilibrio simétrico con muchas empresas, el modelo se aproxima conceptualmente a la competencia monopolística: cada empresa fija su precio ignorando el impacto sobre el conjunto del mercado, los beneficios se anulan por entrada, el precio excede el coste marginal pero iguala el coste medio, y la interacción estratégica directa se diluye. Aun así, el mecanismo competitivo sigue siendo local y espacial, no agregado como en Dixit–Stiglitz. Por eso Salop es mejor entendido como un modelo de competencia monopolística “con microfundamentos espaciales”, no como competencia monopolística pura.
}

El modelo original de diferenciación espacial se debe a HOTELLING.\footnote{%
    En su artículo \textit{Stability in Competition} (1929), modelo que también ha sido empleado en otros campos de la economía como la economía pública (recuérdese el modelo HOTELLING-DOWNS) y la geografía económica.
} 

Este modelo parte de un segmento lineal cuya longitud normalizamos a $1$, en el que se distribuyen uniformemente los consumidores del mercado. La localización de un consumidor representa las preferencias ideales que este tiene respecto a las características del bien. Por tanto, si este debe desplazarse a otro lugar para obtenerlo, debe soportar un coste de transporte por la distancia recorrida: distanciarse de sus características preferidas. Coste que puede representarse mediante la función $t(x)$, donde $x$ es la distancia que recorre el consumidor.\footnote{%
    En general, el modelizar una distribución de preferencias consiste en suponer que el parámetro $x$ sobre el que están definidas las preferencias (la localización en el espacio de características) se distribuye en la economía siguiendo una función de densidad $f(x)$ con una función de distribución $F(x)\in[0,\infty)$, donde $F(0)=0$ y $F(inf)=1$. Por lo. tanto, $F(x)$ es la fracción de consumidores con un parámetro de gustos inferior a $x$
} 

Por simplicidad, supondremos que cada posición del segmento representa una ponderación diferente entre una caracteristica del producto y el resto de sus posibles atributos. De hecho, pensando en el mercado de telefonía podemos suponer que esta característica consiste en la duración de la batería y el resto de atributos normalizados en un componente (por ejemplo, $U_i(\mathbf x, \;\text{resto})$). Supongamos que existen únicamente dos empresas, cuya estructura de costes es idéntica y constante, independientemente del atributo.

Entonces, podemos representar el mercado de la característica como:

\textbf{INTRODUCIR SEGMENTO LINEAL}

Situándose las empresas $a$ y $b$ dentro del espacio: $0\leq a\leq 1-b\leq 1$.

\subsubsection{Costes de transporte: lineales}
El coste total que enfrentará el consumidor por cada bien es el precio que fija la empresa para su producto más el coste de transporte (desutilidad), y la utilidad del consumidor vendrá dada por la diferencia entre su excedente (bruto) y el coste total (precio más costes de transporte):

\eqblock{
\begin{aligned}
    &\begin{aligned}
        &\hspace{2em}T_1(x)=p_1+t(a-x)\quad \text{y}\quad T_1(x)=p_1+t(x-a)\\[1em]
        &T_2(x)=p_2+t(1-b-x)\quad \text{y}\quad T_2(x)=p_2+t(x-(1-b))
    \end{aligned}\\[3em]
    &\text{Utilidad}:\qquad U(x)=\bar U-T_i(x)\qquad \text{donde,}\;
    \left\{\begin{aligned}
        &\bar U\geq\{p_1,p_2\} &&(\substack{\text{\scriptsize si no, no tendría}\\\text{\scriptsize sentido económico}})\\[1em]
        &\text{pero,}\quad \bar U\geq \{T_1(x),T_2(x)\},\;\;\forall x &&(\substack{\text{\scriptsize si no, habría mercado}\\\text{\scriptsize no cubierto}})\\[1em]
    \end{aligned}\right.
\end{aligned}
}{Utilidad del consumidor: Costes lineales y Utilidad de reserva. Modelos espaciales: Ciudad lineal}

Por tanto, la demanda cautiva de cada empresa queda delimitada por el consumidor indiferente entre ambos productos. Ello equivale a buscar la frontera entre las áreas de mercado de las dos empresas para poder luego determinar un equilibrio como resultado de un juego secuencial en dos etapas:
\begin{lnum}
    \item \textbf{Localización}. En primer lugar, la decisión de localización de la empresa; y,
    \item \textbf{Precios}. En segundo lugar, la competencia en precios resultará en el vector de precios de equilibrio. 
\end{lnum}

\textbf{INCLUIR EQUILIBRIO DE CIUDAD LINEAL: COSTES LINEALES}

En el modelo de ciudad lineal con costes de transporte lineales y competencia en precios, no existe equilibrio de Nash en estrategias puras en precios, debido a discontinuidades en la demanda. ¿Por qué? Porque las discontinuidades en precios de la demanda hacen que pequeños cambios en precios provoquen saltos discretos en la cuota de mercado, lo que genera incentivos continuos a subcotizar (undercutting), impidiendo la existencia de un precio óptimo bien definido.\footnote{%
    Si bien si existe uno en estrategias mixtas. Véase: DASGUPTA y MASKIN (1986): \textit{The existence of equilibrium in diconstinous economic games, II: Applications}. The Review of Economic Studies 53, pags. 27-52
} De hecho, existen dos fuerzas contrapuestas, una que empuja hacia la máxima diferenciación (reducción de competencia en precios), y otra que empuja hacia la mínima diferenciación (captación de mercado).

Si bien es cierto que podemos identificar dos resultados intuitivos de referencia o benchmark de poder de mercado:
\begin{la}
    \item \textbf{Principio de máxima diferenciación}. Si cada empresa se situase en un extremo del segmento, el grado de diferenciación sería máximo, lo que maximiza el poder de mercado local y reduce la intensidad de la competencia en precios. Cada empresa actúa casi como monopolista en su mitad del mercado; y,
    \item \textbf{Principio de mínima diferenciación}. El principio de mínima diferenciación implica que ambas empresas se localizan en el centro, que surge por incentivos estratégicos de captación de demanda. No obstante, aunque la diferenciación sea mínima, las empresas siguen teniendo poder de mercado.
\end{la}

\subsubsection{Costes de transporte: Cuadráticos}
Si consideramos costes de transporte no-lineales, por ejemplo cuadráticos, tal que las funciones de desutilidad presentan la forma:\footnote{Modelo inicialmente formulado por D’ASPREMONT, GABSZEWICZ y THISSE (1979)
}

\eqblock{
\begin{aligned}
    &\begin{aligned}
        &\hspace{2em}T_1(x)=p_1+t\cdot (a-x)^2\quad \text{y}\quad T_1(x)=p_1+t\cdot (x-a)^2\\[1em]
        &T_2(x)=p_2+t\cdot (1-b-x)^2\quad \text{y}\quad T_2(x)=p_2+t\cdot (x-(1-b))^2
    \end{aligned}\\[3em]
    &\text{Utilidad}:\qquad U(x)=\bar U-T_i(x)\qquad \text{donde,}\;
    \begin{aligned}
        &\bar U\geq\{p_1,p_2\} &&(\substack{\text{\scriptsize si no, no tendría}\\\text{\scriptsize sentido económico}})\\[1em]
        &\text{pero,}\quad \bar U\geq \{T_1(x),T_2(x)\},\;\;\forall x &&(\substack{\text{\scriptsize si no, habría mercado}\\\text{\scriptsize no cubierto}})\\[1em]
    \end{aligned}
\end{aligned}
}{Utilidad del consumidor: Costes lineales y Utilidad de reserva. Modelos espaciales: Ciudad lineal}

Se nos ofrece una salida a la inexistencia de Equilibrio de NASH por estrategias puras.\footnote{%
    La resolución del problema se pasa a los anexos: [\authorfont{Ver Anexo \ref{anx:modelos_espaciales}}]
} Con costes de transporte cuadráticos, la estructura de la demanda pasa a ser suave y bien comportada respecto a precios y localizaciones. A diferencia del caso con costes lineales, donde pequeñas variaciones en precios provocan saltos discretos en la demanda, los costes cuadráticos hacen que la demanda de cada empresa sea continuamente diferenciable tanto en precios como en localización. Esto elimina el incentivo a subcotizar infinitesimalmente para robar todo el mercado, porque la demanda responde de manera gradual. 

Por tanto, dado este entorno regular, el problema de localización de cada empresa presenta un trade-off bien definido:
\begin{la}
    \item \textbf{Efecto demanda}. Si una empresa se mueve hacia el centro del segmento, se acerca al consumidor indiferente y reduce la distancia media a los consumidores, lo que incrementa su demanda para precios dados: este es el efecto demanda; pero,
    \item \textbf{Efecto estratégico}. Al moverse hacia el centro también reduce la diferenciación respecto a su rival, intensificando la competencia en precios. Esto induce al rival a bajar su precio, lo que reduce la demanda y el margen de la empresa que se ha movido: este es el efecto estratégico.
\end{la} 

Con costes cuadráticos, ambos efectos son continuos y comparables, lo que permite analizar su peso relativo de forma precisa: (i) si prevalece el primer efecto, conduce a una estrategia de producción en masa; y, (ii) si prevalece el segundo efecto, se perseguirá una estrategia de nicho de mercado. 

\textbf{INCLUIR GRÁFICO RESOLUCIÓN: COSTES CUADRÁTICOS}

El modelo de costes cuadráticos se estructura de forma tal que el segundo efecto prevalece, por lo que el equilibrio se da con máxima diferenciación:\footnote{%
    Si prevaleciese el efecto contrario, el resultado recibiría el nombre de principio de diferenciación mínima. Por ejemplo, se ha empleado para explicar por qué los partidos políticos mayoritarios a menudo se parecen mucho entre sí: tienden a situarse en el centro del espectro político con el fin de atraer al mayor número de votantes.
} cada empresa tiene incentivos unilaterales a alejarse del centro y moverse hacia el extremo más cercano, dado el emplazamiento de la otra. Como ambas empresas razonan de la misma forma, el único perfil de estrategias estable es aquel en el que cada empresa se sitúa en un extremo del segmento, dando lugar a diferenciación máxima.\footnote{%
    La existencia del equilibrio puro se asegura, por tanto, por dos razones fundamentales. Conceptualmente, porque los costes cuadráticos suavizan la competencia y eliminan las discontinuidades que impedían la estabilidad estratégica. Analíticamente, porque la función de beneficios es continua y cuasi-cóncava en las estrategias de localización, lo que garantiza la existencia de un máximo interior o de esquina bien definido. El equilibrio en localización con diferenciación máxima no es un resultado ad hoc, sino la consecuencia de que las empresas internalizan cómo su posición afecta estratégicamente al comportamiento en precios de su rival. En este sentido, el modelo muestra con claridad que la diferenciación no es solo una respuesta a preferencias heterogéneas, sino un instrumento estratégico para relajar la competencia en precios.
}

\subsubsection{Implicaciones de Política Económica}
Por tanto, desde un punto de vista social, en ninguno de los casos coincide el equilibrio de mercado con el óptimo social. Con costes lineales, la competencia empuja hacia mínima diferenciación (excesiva aglomeración), mientras que con costes cuadráticos el equilibrio privado lleva a diferenciación máxima (excesiva dispersión). En ambos casos, la localización en equilibrio responde a incentivos estratégicos (captación de demanda o relajación de la competencia en precios), no a la eficiencia asignativa. Esto revela una distorsión espacial propia de la competencia monopolística: la diferenciación elegida por las empresas no internaliza los costes de transporte sociales, justificando la intervención.

El planificador social, internalizando los costes de transporte soportados por los consumidores, elegirá las localizaciones que minimizan los costes totales entre consumidores y empresas. Colocar a cada empresa en el centro de su mitad reduce la suma de costes de transporte agregados, independientemente de si estos costes son lineales o cuadráticos; lo que cambia entre ambos casos es la penalización de la distancia, no la localización óptima que minimiza el coste agregado.

\textbf{INCLUIR GRÁFICO ÓPTIMO SOCIAL: COSTES LINEALES Y CUADRÁTICOS}

\paragraph{Limitaciones}
No obstante, el rasgo estructural del modelo es la presencia de fronteras. Esta característica aparentemente inocua introduce un problema profundo: las empresas no enfrentan la misma presión competitiva en todas las localizaciones ya que cada una compite solo por un lado. Analíticamente, esto implica que la demanda residual y su elasticidad dependen no solo del precio, sino también de la posición en el segmento. Como consecuencia, los incentivos a fijar precios, a desplazarse espacialmente o a entrar en el mercado no son simétricos entre empresas idénticas. 

Por ejemplo, una empresa puede preferir localizarse cerca de un extremo no porque los consumidores estén allí, sino porque la frontera actúa como una protección natural frente a la competencia. En términos analíticos, la frontera reduce la elasticidad de la demanda residual y permite márgenes más altos. Así, los resultados sobre mínima o máxima diferenciación dejan de ser una consecuencia general de las preferencias o de la tecnología, y pasan a depender de las especificaciones del modelo, lo que explica por qué el mismo modelo puede predecir aglomeración en el centro o dispersión en los extremos bajo supuestos ligeramente distintos.

Estas propiedades limitan seriamente la utilidad del enfoque lineal para analizar mercados con muchas variedades y entrada libre. Cuando hay más de dos empresas, las de los extremos disfrutan sistemáticamente de una ventaja competitiva que no tiene un análogo realista mercados de bienes diferenciados.

\subsection{Número de empresas endógeno: Modelo ciudad circular}
Precisamente para superar estas limitaciones surge el modelo de Ciudad Circular desarrollado por SALOP en su artículo \textit{Monopolistic Competition with Outside Goods} de 1979. La motivación es eliminar los artefactos de borde y garantizar que cada empresa enfrente una competencia local simétrica. De este modo, se puede estudiar la competencia monopolística con entrada libre y muchas variedades, separando claramente el papel de la diferenciación (vía costes de transporte) del de la estructura competitiva. 

Sobre un segmento circular de perimetro igual a uno, sin alterar los supuestos relativos al consumidor del modelo anterior, en relación a la oferta del mercado suponemos que: (i) todas las empresas tienen acceso a la misma tecnología y enfrentan la misma estructura de costes; (ii) se sitúan alrededor del círculo de forma equidistante (por lo tanto, se impone de forma exógena la diferenciación máxima)\footnote{%
    En una ciudad circular, ninguna posición es superior a las demás: no hay extremos ni centro. Al contrario que la ciudad linear donde en el centro existe una mayor demanda, aquí el espacio de los productos no privilegia ninguna posición en particular. Por tanto, aquí no existe el efecto mayor demanda por localización. Ello lleva a que las empresas se sitúen lo más lejos posible de sus competidores para suavizar la competencia en precios. Si hay n empresas, la distancia entre empresas será 1/n, asumiendo que el perímetro del círculo es 1. Esto lo hemos puesto como supuesto, pero en verdad es un resultado del modelo, ya que por el propio interés de las empresas, se colocarán de esta manera. En aras de un mayor realismo, sería preferible que las empresas eligiesen su ubicación, ya sea a la vez que se entra otras la decisión de entrada. Sin embargo, el interés del modelo de SALOP se centra no en la dirección concreta de productos, sino en el estudio del alcance de la entrada. La omisión de la elección de la ubicación permite estudiar la decisión de entrada de una forma simple y tratable. En cambio, la elección de la localización se estudia en el marco del modelo de la ciudad lineal.
}; (iii) existe un número elevado de empresas potenciales, lo que da lugar al supuesto de libre entrada en el mercado.

Como ocurría anteriormente, el modelo es asimilable a un juego con dos periodos: 
\begin{lnum}
    \item En el primer periodo, un número elevado de empresas idénticas deciden si entran en el mercado, entrando finalmente $n$ de ellas que se sitúan de forma equidistante unas de otras en el círculo; y,
    \item En el segundo periodo, tras observar la localización respectiva, se compite en precios.
\end{lnum}

La resolución del juego emplea la inducción hacia atrás.\footnote{%
    La solución analítica se encuentra en [\authorfont{Ver Anexo \ref{anx:modelos_espaciales}}]
} Desde un punto de vista positivo, en la ciudad circular de Salop cada empresa elige precio sabiendo que compite localmente con sus dos vecinas: los consumidores compran una unidad a la firma que minimiza sus costes totales a lo largo del círculo. En este entorno, el precio en equilibrio es un mark-up sobre el coste marginal que depende: (i) de la intensidad de la competencia local; y, (ii) del grado de diferenciación espacial inducido por la distancia media a la variedad más cercana. 

En este caso particular, el margen precio–coste marginal actúa principalmente como transferencia (reduce utilidad neta del consumidor y aumenta beneficios), mientras que el bienestar total se juega sobre todo en la distancia/variedad (costes de transporte) y en los costes fijos de entrada. 

\subsubsection{Implicaciones de Política Económica}
Desde el punto de vista normativo, el planificador elige el número de empresas $n$ para minimizar el coste social, costes de transporte más costes fijos. Con $n$ firmas simétricamente espaciadas, la distancia media al proveedor más cercano cae con el número de empresas $n$, y el coste total de transporte puede escribirse:

\eqblock{
\begin{aligned}
    T=n\cdot 2\cdot \int_{0}^{(1/2)\cdot n}t\cdot x\cdot \mathrm dx=n\cdot t\left(\frac{1}{2n}\right)^{2}=\frac{t}{4n}
\end{aligned}
}{Costes de transporte de los consumidores. Modelos espaciales: Ciudad circular (SALOP)}

El óptimo social aparece donde el beneficio marginal de una firma adicional iguala su coste fijo marginal. Con libre entrada, la decisión privada de entrar no internaliza dos externalidades opuestas, existiendo dos efectos que inciden sobre el bienestar social en sentido opuesto: 
\begin{la}
    \item \textbf{Inapropiabilidad del excedente del consumidor}. Cuando una empresa considera la entrada al mercado no internaliza el beneficio generado a través de la externalidad positiva sobre los consumidores (mayor variedad), lo que empuja hacia menos entrada; y,
    \item \textbf{El robo de clientes}. Una externalidad negativa sobre las incumbentes, ya que el entrante captura demanda que antes era de otros y reduce beneficios del incumbente, lo que empuja hacia demasiada entrada porque el entrante ignora esa pérdida ajena.
\end{la}

En el modelo circular canónico de SALOP, el efecto de robo de negocio tiende a dominar bajo los supuestos estándar, de modo que el equilibrio competitivo con entrada libre suele exhibir entrada excesiva respecto al planificador.
 














\clearpage
\section{Conclusión} 

\subsection{Extensiones y relación con otras partes del Temario}


\subsection{Opinión}



\subsubsection{Idea final (Salida o cierra)}


\clearpage
\pagenumbering{gobble}
\MyBibliography
\MyBibItem{Autor}{Título}{Año}{DOI -- LINK}




%ANEXOS
\clearpage
\anexo{\textbf{Preguntas Test}}
%\input{\TemaCode/questions.tex} 


\clearpage
\anexo{Formulación General: Modelo DIXIT-STIGLITZ}



\textbf{Esto de aquí abajo es lo que tenemos que introducir bien. Es el modelo baseline con preferencias por la variedad}

Por otro lado, en la literatura, muchos autores han impuesto las tres restricciones juntas en los que NEARY (2004) llama “Dixit-Stiglitz lite”. En cuyo caso tendríamos la especificación siuiente para la función de utilidad:
 
La función de tipo CES verifica, igual que antes, la preferencia por la variedad (la utilidad que le reporta al individuo el consumo de n variedades es, ceteris paribus, superior a la que le reporta el consumo de tan sólo n – 1 variedades)  : por lo que los consumidores desearán consumir todas las variedades. 
De este modo:
0<(σ-1)/σ<1⟹1<σ
para que los bienes sean sustitutivos imperfectos, consecuencia del supuesto de curvas de indiferencia convexas al origen [ver Tema 3.A.8] y, por tanto, de una función de utilidad cóncava.  Es más, sabemos que: ρ=((σ-1))⁄σ va a indicar el grado de preferencia porla variedad. Esto es:
	Si ρ⟶0 (σ⟶1): el agente muestra una muy elevada preferencia por la variedad; y,
	Si ρ⟶1 (σ⟶∞):  los bienes están cerca de ser percibidos como sustitutos perfectos por el agente.

Dicho esto, desarrollaremos aquí la versión general (por ser la que se desarrolla en los apuntes de CECO) y dejamos para el Anexo las consecuencias de considerar el otro tipo de definición de la función de utilidad.


\clearpage
\anexo{Índice de Precios: El Deflactor del Bienestar en el Modelo DIXIT-STIGLITZ}\label{anx_deflactor_bienestar}
La expresión “deflactor del bienestar” suena técnica, pero en el modelo Dixit-Stiglitz nombra una idea bastante precisa: el índice de precios no se interpreta simplemente como una media de precios de mercado, sino como el objeto que transforma magnitudes nominales en magnitudes reales de utilidad o, dicho de otro modo, como el coste mínimo de obtener una determinada unidad del agregado de consumo. Por eso “deflacta” el bienestar: sirve para pasar del ingreso monetario observado al poder de compra relevante en términos de satisfacción.

El punto de partida es que, en este tipo de modelos, al consumidor no le importa solo cuánto gasta, ni siquiera solo cuántas unidades físicas adquiere, sino la utilidad que obtiene del conjunto de variedades. Eso obliga a distinguir entre dos cosas que en modelos más simples casi coinciden: el valor monetario del consumo y la cantidad “real” de bienestar que ese gasto permite alcanzar. Si todos los bienes fueran idénticos y hubiera un único precio, esta distinción sería trivial. Pero cuando el consumo se compone de muchas variedades diferenciadas, con una agregación CES, la pregunta correcta deja de ser “¿cuánto consume?” y pasa a ser “¿cuánto cuesta alcanzar un cierto nivel del agregado de consumo que entra en la utilidad?”.

Eso es exactamente lo que responde el índice de precios CES. No aparece por conveniencia algebraica, sino como solución de un problema dual muy importante. En lugar de formular el problema habitual del consumidor —maximizar utilidad sujeta a una restricción presupuestaria—, puede plantearse el problema dual: minimizar el gasto necesario para alcanzar un nivel dado de utilidad o de subutilidad. En el bloque de variedades, este problema consiste en elegir las cantidades de cada variedad de manera que se alcance un nivel prefijado del agregado CES al menor coste posible. El resultado de esa minimización es una función de gasto. Y esa función de gasto adopta una forma extremadamente simple: gasto mínimo igual al índice de precios multiplicado por la cantidad del agregado. En símbolos, si llamamos C al agregado CES de variedades, el gasto mínimo para alcanzar C es E(P,C)=P\,C. Por eso el índice P puede interpretarse como el “precio unitario” del bienestar dentro del bloque de variedades: es lo que cuesta comprar una unidad del compuesto de consumo relevante para la utilidad.

Aquí está el corazón de la idea. El bienestar, en este modelo, no se identifica con el dinero gastado, sino con la cantidad del agregado C que el consumidor puede permitirse. Si el gasto mínimo necesario para obtener una unidad de C es P, entonces, dado un gasto total E, la cantidad del agregado que puede comprarse es C=E/P. Esto no es una aproximación intuitiva; es una identidad exacta derivada de la función de gasto dual. Por eso se dice que P “deflacta” el bienestar: divide una magnitud monetaria, E, para convertirla en una magnitud real, C, que sí tiene significado directo en la utilidad. Cuando además el bloque CES es el único bien, o cuando el gasto destinado a él ya está determinado en una primera etapa, entonces la utilidad relevante del consumidor crece con E/P. En ese sentido, P funciona como un deflactor del ingreso o del gasto en términos de bienestar.

Se entiende mejor si se compara con el uso ordinario de un deflactor en macroeconomía. Cuando se deflacta el PIB nominal por un índice de precios, lo que se intenta medir es cuánto producto “real” representa una cantidad monetaria dada. Aquí sucede algo análogo, pero en vez de transformar dinero en “cantidad física de producción”, el índice CES transforma dinero en “capacidad de alcanzar utilidad del bloque de variedades”. Por eso es un deflactor del bienestar. No mide directamente la competencia; mide cuán caro o cuán barato es convertir renta nominal en satisfacción real.

Ahora bien, ¿por qué el índice de precios puede desempeñar ese papel si no es simplemente una media aritmética de precios? Precisamente porque la agregación CES incorpora el patrón de sustitución del consumidor. Si algunos bienes se encarecen, el consumidor puede reorientar parte de su gasto hacia otros. Si aparecen nuevas variedades, el consumidor dispone de más combinaciones para alcanzar el mismo nivel del agregado. El índice P resume toda esa información en una sola magnitud: no es el precio de una variedad concreta, ni tampoco la intensidad competitiva entre empresas, sino el coste mínimo de conseguir una unidad de la cesta “ideal” implícita en las preferencias. Dicho de manera rigurosa, el índice de precios es la función de coste unitario asociada al agregador de cantidades. Por eso es el objeto correcto para medir el precio del bienestar, mientras que un precio individual no lo sería.

Con esto ya puede verse por qué una disminución de P implica una mejora del bienestar y un aumento de P un empeoramiento, todo lo demás constante. Si el ingreso monetario o el gasto nominal en variedades se mantiene fijo, una caída de P significa que el mismo dinero permite alcanzar una cantidad mayor del agregado C, porque C=E/P. Y como la utilidad es creciente en C, el bienestar aumenta. Inversamente, si P sube, el mismo gasto permite comprar menos del agregado CES, luego disminuye la utilidad alcanzable. No es una afirmación verbal, sino una consecuencia exacta de la relación dual entre utilidad y gasto. La caída del índice de precios no tiene por qué significar que cada empresa individual haya reducido su precio nominal; puede significar, por ejemplo, que la combinación total de precios y variedad hace más barato alcanzar una unidad de utilidad. Y esa es justamente la razón por la que el bienestar mejora.

Aquí conviene insistir en algo que suele generar confusión: en Dixit-Stiglitz, el índice de precios puede disminuir aunque la intensidad competitiva, en el sentido de markup sobre coste marginal, permanezca constante. Eso ocurre porque el índice recoge no solo los precios individuales, sino también la estructura del conjunto de variedades disponibles. Si aumenta el número de variedades y los precios individuales no cambian, el consumidor tiene más opciones para componer su cesta. Bajo preferencias CES, esa mayor variedad reduce el coste mínimo de alcanzar una unidad del agregado. Por eso P puede caer sin que haya habido una “guerra de precios”. Lo que se abarata no es necesariamente cada bien por separado, sino el acceso a un cierto nivel de bienestar. Esto es fundamental: el índice de precios es un precio del agregado de utilidad, no un termómetro puro de rivalidad estratégica.

Puede verse con especial claridad en el caso simétrico, cuando todas las variedades tienen el mismo precio p y existen N variedades. Entonces el índice CES adopta la forma P=p\,N^{1/(1-\sigma)}. Como \sigma>1, el exponente 1/(1-\sigma) es negativo, de modo que al aumentar N, el término N^{1/(1-\sigma)} disminuye. Si p permanece constante, también disminuye P. ¿Qué significa esto? No que cada empresa esté cobrando menos por su producto, sino que el coste de alcanzar una unidad del agregado de consumo es menor porque el consumidor puede distribuir su gasto entre más variedades y eso, dadas las preferencias, genera más utilidad. En otras palabras, el mismo euro “rinde más” en términos de bienestar. Eso es exactamente lo que quiere decir que P sea un deflactor del bienestar.

A estas alturas quizá se vea mejor por qué no debe confundirse el índice de precios con la competencia. La competencia, en este modelo, viene resumida por la elasticidad de sustitución \sigma, que determina el markup de equilibrio. Si \sigma es alta, las variedades son buenas sustitutas, la demanda percibida por cada empresa es muy elástica y el margen sobre coste marginal es pequeño. Si \sigma es baja, las variedades son menos sustituibles, la demanda es menos elástica y el margen es mayor. Ese es el margen competitivo. Pero el bienestar del consumidor no depende solo de si las empresas tienen mucho o poco poder de mercado. También depende de cuántas variedades existen y de cuánto cuesta, en conjunto, acceder al agregado de consumo que esas variedades forman. El índice P resume esta segunda dimensión. Por eso una caída de P es buena para el consumidor aunque no haya cambiado la intensidad de la competencia: el poder de compra relevante para la utilidad ha aumentado.

También ayuda pensar en la relación entre utilidad indirecta y función de gasto. En teoría del consumidor, la utilidad indirecta expresa el nivel máximo de utilidad alcanzable dados precios e ingreso. La función de gasto expresa el gasto mínimo necesario para alcanzar una utilidad dada. Son dos caras de la misma moneda. En preferencias homotéticas como las CES, esta dualidad es especialmente limpia. El índice de precios aparece como la función de coste unitario, y eso hace que la utilidad indirecta dependa del cociente entre ingreso e índice de precios. Así, si la utilidad indirecta se escribe, esencialmente, como una función creciente de Y/P, entonces P es el deflactor natural del bienestar: cuanto menor es, mayor es la utilidad máxima alcanzable con el mismo ingreso nominal. En este sentido, decir que “baja el índice de precios” equivale a decir que “sube el ingreso real relevante para la utilidad”.

Quizá una fuente adicional de desconcierto sea que, en lenguaje corriente, “precio” sugiere siempre el importe que se paga por una unidad física observable. Aquí no. Aquí el “precio” del índice es el precio de una unidad abstracta del agregado CES, una unidad de capacidad de consumo relevante para la utilidad. Por eso puede moverse por cambios en la variedad aunque ningún precio individual cambie. No es el precio de un bien concreto; es el precio sombra de una unidad de bienestar dentro del bloque de bienes diferenciados. Una vez se entiende esto, la lógica deja de parecer paradójica: el bienestar mejora cuando el mismo ingreso compra más de ese agregado y empeora cuando compra menos.

Dicho de la forma más rigurosa posible, la razón última por la que P mide bienestar es que está definido dualmente a partir de las preferencias. No es un índice elegido arbitrariamente para resumir datos; es la magnitud exacta que sale del problema de minimización del gasto compatible con un nivel dado del agregado de utilidad. Y la razón por la que su descenso mejora el bienestar es que reduce el coste unitario de ese agregado. Si el coste de una unidad de utilidad baja, entonces, con el mismo ingreso, puede alcanzarse más utilidad. Si ese coste sube, sucede lo contrario. Todo el contenido económico del “deflactor del bienestar” está ahí.

La frase más compacta y exacta sería esta: en Dixit-Stiglitz, el índice de precios es el deflactor del bienestar porque representa el coste mínimo de comprar una unidad del agregado de consumo que entra en la utilidad; en consecuencia, el bienestar real viene dado por el gasto nominal dividido por ese índice, de modo que una caída de P aumenta el poder adquisitivo en términos de utilidad y una subida lo reduce.

Si te sirve, en el siguiente paso puedo derivártelo formalmente, desde el problema de minimización del gasto, hasta llegar a E(P,C)=PC y a la utilidad indirecta, sin saltos.



\clearpage
\anexo{Elasticidad de sustitución: ¿Qué es? ¿Qué implicaciones tiene? -- Modelo DIXIT-STIGLITZ}

\textcolor{red}{\textbf{OJO!}} Aquí sería muy importante desarrollar el concepto de elasticidad de sustitución entre variedades, las implicaciones que tiene y por qué se escoge una elasticidad tipo CES como referencia al desarrollar el modelo: diferencias intersectoriales? diferencias de consumo?

















\clearpage
\anexo{Modelos espaciales: Resolución analítica}\label{anx:modelos_espaciales}
En este anexo se formula la resolución analítica de los modelos espaciales: HOTELING y SALOP.


\textsc{Modelo de ciudad lineal: HOTELING}

Supongamos: (i) un continuo de consumidores se distribuyen de forma uniforme en un intervalo linear $[0,1]$;\footnote{Se asume distancia igual a uno por normalizar, por sencillez analítica.
} (ii) cada punto del intervalo representa una variedad de un bien que refleja la combinación de dos características (siguiendo la aproximación de la teoría de la demanda de características de LANCASTER);\footnote{%
    En relación a este supuesto, una crítica al modelo es la unidimensionalidad de la diferenciación. En este modelo las empresas compiten en un segmento. En línea con esto, el modelo de DOWNS-HOTELLING, que es una aplicación a la política de este modelo, supone que los partidos políticos compiten en un segmento de izquierda a derecha, pero se podría argumentar que existen distintos segmentos (p.ej. en el ámbito económico puedes ser más liberal o más intervencionista y en el ámbito social puedes ser más progresista o más conservador) [ver temas 3.A.24 y 4.B.2].
} (iii) las variedades efectivas se encuentran situadas en la localización del segmento donde decidan situarse las empresas; (iv) cada consumidor dispone de un excedente bruto (utilidad reserva) y unos costes de transporte asociados a la distancia que deben recorrer entre su localización y la variedad ofrecida por la empresa más cercana (conceptualmente, \textit{lo que me cuesta distanciarme de mi variedad ideal}); y, (v) existen dos empresas con igual tecnología de producción.

Existen difentes formulaciones del modelo aunque las más relaventes son:
\begin{la}
    \item \textbf{Localización exógena}. Con localización exógena, podemos llegar a un equilibrio tanto con costes cuadráticos como con costes lineales; y,
    \item \textbf{Localización endógena}. Con localización endógena, el equilibrio no queda garantizado para costes lineales ya que la demanda es discontinua. Por tanto, la única alternativa viable en términos analíticos es los costes cuadráticos.
\end{la}

Desarrollaremos únicamente el equilibrio para costes cuadráticos y localización endógena. Se distinguen dos etapas en el problema:
\begin{lnum}
    \item \textbf{Decisión de localización}. En esta primera etapa, la empresa decide donde se situará en el espacio de variedades ($a,b$); y,
    \item \textbf{Decisión de precios}. En esta segunda etapa la empresa elige su precio tal y que maximice los beneficios. Dadas ($a,b$) las empresas deciden ($p_A,p_B$)
\end{lnum}

Por inducción hacia atrás, la etapa relevante para entender los incentivos de localización es la segunda, porque los precios reaccionan endógenamente a los cambios en $a,b$. 

Dadas las localizaciones, el consumidor indiferente $\tilde x$ se define por igualdad de precios totales:

$$\begin{aligned}
    &\boxed{p_A+t(\tilde x-a)^2=p_B+t(\tilde x -b)^2}\\[1em]
    &\overset{\substack{\text{\scriptsize Resolviendo, obtenemos}\\\text{\scriptsize la frontera del mercado}}}{\Longrightarrow}\left\{\begin{aligned}
        &D_A=\tilde x:&&\frac{p_B-p_A-t\cdot a^2+t\cdot b^2}{2\cdot t\cdot (b-a)}\\
        &D_B=1-\tilde x: &&\frac{p_A-p_B-t\cdot b^2+t\cdot a^2}{2\cdot t\cdot (a-b)}\\
    \end{aligned}\right.\\[2em]
    &\text{Entonces,}\qquad \max \Pi_i(a)=(p_i-c)\cdot D_i(p_i,p_j,a,b)\;\overset{\text{\scriptsize CPO's}}{\Longrightarrow}\;\frac{\partial\Pi_i}{\partial p_i}=0\\[1em]
    &\substack{\text{En equilibrio de precios,}\\\text{ los precios son funciones óptimas}\\\text{ de las localizaciones}}\;\Longrightarrow\;p_i=p_i(a,b)
\end{aligned}$$
 
Ahora, en la primera etapa, la empresa $A$ elige a anticipando cómo los precios reaccionan:

$$\begin{aligned}
    \underset{\text{\scriptsize Análogamente para $B$}}{\text{Su beneficio relevante es:}}\qquad \Pi_A(a)=\Pi_A(p_A(a,b),p_B(a,b), a, b)\\[1em]
    \overset{\text{\scriptsize CPO's}}{\Longrightarrow} 
    \frac{\partial \Pi_A}{\partial a}=\qquad \underbrace{\underbrace{\frac{\partial \Pi_A}{\partial p_A}}_{\substack{\text{\scriptsize Como $p_A$ es óptimo en la}\\\text{\scriptsize segunda (primera) etapa: $\frac{\partial \Pi_A}{\partial p_A}=0$}}}\frac{\partial p_A}{\partial a}}_{=0}+\underbrace{\frac{\partial \Pi_A}{\partial p_B}}_{\frac{\partial \Pi_A}{\partial p_B}=\frac{\partial \Pi_A}{\partial D_A}\frac{\partial D_A}{\partial p_B}}\frac{\partial p_B}{\partial a}+\frac{\partial \Pi_A}{\partial D_A}\frac{\partial D_A}{\partial a}\\
    \underset{\text{\scriptsize Sustituyendo}}{\Longrightarrow}\frac{\partial \Pi_A}{\partial a}=\frac{\partial \Pi_A}{\partial D_A}\cdot \left(\frac{\partial D_A}{\partial p_B}\frac{\partial p_B}{\partial a}+\frac{\partial D_A}{\partial a}\right)
\end{aligned}$$

Por tanto, si una empresa decida relocalizarse, experimentará dos efectos contrapuestos:\footnote{%
    Si los precios son fijados exógenamente, iguales a $\bar p$ (por ejemplo, debido a regulación), prevalece el principio de mínima diferenciación, ya que solo existe el efecto demanda. Esto es relevante para el test. Si en un test preguntan por el equilibrio con precios fijos, el equilibrio sería las dos empresas en el centro. Esto es así, porque en este caso solo estaría presente el efecto demanda (o efecto cuota de mercado), y no el efecto estratégico que es el que causa un incentivo para que las empresas se diferencien más.
}
\begin{la}
    \item Efecto demanda: $\frac{\partial D_A}{\partial a}$. Captura el efecto directo de mover la localización sobre la frontera $\tilde x$. Si A se mueve hacia el centro, reduce la distancia media a los consumidores cercanos y, con costes cuadráticos, ese efecto será positivo: aumentará la demanda de A. Dicho de otra forma, para precios dados, cada empresa querría acercarse al centro para captar consumidores lo que, con costes lineales, equivaldría al principio de mínima diferenciación; y,
    \item Efecto estratégico: $\frac{\partial D_A}{\partial p_B}\frac{\partial p_B}{\partial a}$. Que recoge el efecto indirecto: si A se acerca a B, los productos sevuelven menos diferenciados, lo que hace que las demandas de las empresas se vuelvan más sensibles al precio. Dicho de otra forma, al acercarse a su rival, le fuerza a competir más agresivamente en precios, lo que reduce el beneficio de la empresa. 
\end{la}

D’ASPREMONT et al. (1979) muestran que, para costes de transporte cuadráticos, el equilibrio es tal que las dos empresas se sitúan en los dos extremos de la ciudad (diferenciación máxima). Es decir, 

$$\begin{aligned}
    &\text{ Como, $\frac{\partial \Pi_A}{\partial D_A}>\underbrace{p_A-c}_{\text{\scriptsize Mark-up positivo}}>0$}\;\Longrightarrow\;\substack{\text{El signo del incentivo depende exclusivamente}\\\text{ del término entre paréntesis}}\\[1em]
    &\frac{\partial \Pi_A}{\partial a}=(p_A-c)\cdot \left(\underbrace{\frac{\partial D_A}{\partial a}}_{\text{\scriptsize Efecto demanda}}+\underbrace{\frac{\partial D_A}{\partial p_B}\frac{\partial p_B}{\partial a}}_{\text{\scriptsize Efecto estratégico}}\right)\quad \text{donde,}\;
    \left\{\begin{aligned}
        &\frac{\partial D_A}{\partial a}>0\equiv&&
        \underset{\text{\scriptsize $\Longrightarrow$ Efecto que siempre empuja a mínima diferenciación}}{\left(\begin{aligned}
            &\text{\scriptsize Si A se mueve hacia el centro,}\\
            &\text{\scriptsize se acerca a más consumidores}
        \end{aligned}\right)}\\[1em]
        &\frac{\partial D_A}{\partial p_B}>0\equiv&&
        \underset{\text{\scriptsize $\Longrightarrow$ No hay disuasión en este punto}}{\left(\begin{aligned}
            &\text{\scriptsize Si B sube su precio, algunos consumidores}\\
            &\text{\scriptsize cambian a A ($\sim$ Efecto sustitución)}
        \end{aligned}\right)}\\[1em]
        &\frac{\partial p_B}{\partial a}<0\equiv&&
        \underset{\text{\scriptsize $\Longrightarrow$ B responde bajando su precio}}{\left(\begin{aligned}
            &\text{\scriptsize Si se acerca a B, menos diferenciación}\\
            &\text{\scriptsize lo que hace que demanda de B más elástica}
        \end{aligned}\right)}
    \end{aligned}\right.\\[2em]
    &\text{Por tanto,}\qquad \underbrace{\frac{\partial D_A}{\partial p_B}\frac{\partial p_B}{\partial a}<0}_{\text{\scriptsize Si te acercas a tu rival, compites en precios y esto te perjudica}}\qquad \text{y}\qquad \boxed{\left|\frac{\partial D_A}{\partial p_B}\frac{\partial p_B}{\partial a}\right|>\left|\frac{\partial D_A}{\partial a}\right|\Longrightarrow\;\frac{\partial \Pi_A}{\partial a}<0}
\end{aligned}$$

Por tanto, como vemos, con costes cuadráticos la reducción en diferenciación dispara la competencia en precios y la caída del precio afecta a todas las unidades vendidas, no solo a las marginales. Por eso, el efecto estratégico es de primer orden (muy relevante), mientras que el efecto demanda es más débil. Dicho de otra forma, mientra que el efecto demanda es local (se ganan algunos consumidores marginales), el efecto estratégico es global en precios (afecta al precio óptimo en todo el mercado), lo que conduce a un equilibrio de máxima diferenciación.

En conclusión, en el modelo de ciudad lineal con costes cuadráticos, aunque mover la localización hacia el centro aumenta la demanda directa de la empresa, dicho movimiento reduce la diferenciación y endurece la competencia en precios. La reacción estratégica del rival en precios domina al efecto directo sobre la demanda, de modo que el beneficio disminuye al acercarse al centro. Por ello, las empresas tienen incentivos a maximizar la diferenciación espacial.

\textsc{Modelo ciudad circular: SALOP}

El modelo de ciudad circular, desarrollado por SALOP es similar al de su predecesor, pero introduce dos diferencias principales: un continuo de consumidores se distribuyen de forma uniforme en un intervalo circular de perímetro unitario. SALOP introduce una ciudad circular debido a que la ciudad lineal de HOTELLING no sería apropiada para su estudio debido a los efectos frontera (boundary effects). La circunferencia propuesta por SALOP solventa este problema al no tener fronteras. Igual que en el modelo anterior, la solución consiste en dos etapas que se resuelven mediante inducción hacia atrás. 

En la primera etapa, como todas las empresas están equidistantes, buscamos un equilibrio simétrico donde todos cobren el mismo precio. Para hallar la demanda de una variedad, encontramos al consumidor indiferente ($\tilde x$) entre dos variedades (resolviendo como en el caso de ciudad linear para uno de los espacios entre dos empresas), y multiplicamos por 2 (ya que cada empresa se enfrentará a demanda por ambos lados y de cuantías idénticas por el supuesto de simetría). De este modo, suponiendo que el resto de empresas fijan un precio p39, la empresa i se enfrentará a la siguiente demanda:

$$\begin{aligned}
    &p_i+t\cdot\tilde x=p+t\left(\frac{1}{n}-\tilde x\right)\Longrightarrow t\left(2\tilde x-\frac{1}{n}\right)=p-p_i\Longrightarrow \tilde  x=\frac{1}{2n}+\frac{p-p_i}{2t}\\[1em]
    &\overset{\substack{\text{\scriptsize Como $i$ tiene el}\\\text{\scriptsize mismo \textit{consumidor}}\\\text{\scriptsize por ambos lados}}}{\Longrightarrow}\underset{\text{\scriptsize Exige mercado cubierto $\equiv\;0\leq \tilde x\leq \frac{1}{n}\sim |p_i-p|\leq \frac{t}{n}$}}{\text{Demanda total:}}\qquad D_i=2\cdot \tilde x=\frac{p-p_i}{t}+\frac{1}{n}\\[2em]
    &\text{Primera (segunda) etapa, número de empresas dado ($n$) competencia en precios:}\\
    &\over\begin{aligned}
        \\
        &\max_{p_i}\Pi_i(p_i,p)=(p_i-c)\left(\frac{1}{n}+\frac{p-p_i}{t}\right)=(p_i-c)D_i-F\\[1em]
        &\overset{\text{\scriptsize CPO's}}{\Longrightarrow}\frac{\partial\Pi_i}{\partial p_i}=0\Longrightarrow\frac{\partial \Pi_i}{\partial p_i}=\left(\frac{1}{n}+\frac{p-p_i}{t}\right)+(p_i-c)\left(-\frac{1}{t}\right)=0\\[1em]
        &\overset{\text{\scriptsize Reordenando}}{\Longrightarrow}\frac{1}{n}+\frac{p-p_i}{t}-\frac{p_i-c}{t}=0\Longrightarrow\frac{1}{n}+\frac{p+c+2p_i}{t}=0\;\Longrightarrow\;\boxed{p_i=\frac{p+c}{2}+\frac{t}{2n}}\\[1em]
        &\text{Imponiendo equilibrio simétrico, $p_i=p$}:\qquad p=\frac{p+c}{2}+\frac{t}{2n}\Longrightarrow \boxed{p=c+\frac{t}{n}}\\[1em]
        &\text{Sustituyendo $p(n)$ en $\Pi_i$:}\qquad \Pi_i=(p_i-c)\left(\frac{p-p_i}{t}+\frac{1}{n}\right)-F\Longrightarrow\boxed{\Pi_i=\frac{t}{n^2}-F}\\[1em]
    \end{aligned}
\end{aligned}$$

Por tanto, una vez analizada la competencia en precios dado el número de empresas, debemos determinar la decisión de entrada mediante la condición de beneficio nulo. 

$$\begin{aligned}
    &\text{Segunda (primera) etapa, decisión de entrada:}\\
    &\over\begin{aligned}
        \\
        &\Pi_i(n)=\frac{t}{n^2}-F\qquad \text{y,}\qquad D_i=\frac{1}{n}\;\;(\text{\scriptsize En simetría})\\[1em]
        &\Longrightarrow\;
        \left\{\begin{aligned}
            &n^*=\sqrt{\frac{t}{F}}\\[0.5em]
            &p^*=c+\frac{t}{n^*}=c+\sqrt{t\cdot F}
        \end{aligned}\right.
    \end{aligned}
\end{aligned}$$

De este modo, en el equilibrio se da una situación en la que el precio es mayor que el coste marginal, pero debido a la condición de libre entrada, el precio se iguala al coste medio y el beneficio de las empresas es nulo, de forma similar a lo estudiado en el modelo de DIXIT y STIGLITZ. A medida que disminuye F entran más empresas con un precio más bajo. Si F tiende a cero, en el límite, el mercado se vuelve perfectamente competitivo y todos los consumidores compran su producto preferido. Del mismo modo, a medida que aumenta t, entran más empresas con un precio más alto, ya que aumenta la posibilidad de diferenciación. En cualquier caso, siguen entrando empresas hasta que se cumpla la condición de beneficio nulo.


\end{document}